\chapter{偏振、晶体与界面电磁波}
\section{Jones矢量与偏振椭圆}\label{th:o4jones}
在均匀各向同性介质中，平面波为横波。取传播方向$+z$，固定相位约定$\exp[i(kz-\omega t)]$，省略公共相位后的Jones矢量为
\[
 \vct e=\begin{pmatrix}a\\b\ee^{i\delta}\end{pmatrix},\quad a,b\geq0.
\]
某固定位置以$\omega t$为时间变量，可写$E_x=a\cos\omega t$、$E_y=b\cos(\omega t-\delta)$。Jones矢量只适用于确定的完全偏振分量；自然光的随机偏振不能用一个恒定Jones矢量替代。

若$\delta=0$或$\pi$，两分量同相或反相，末端沿直线运动。若$a=b$且$\delta=\pm\pi/2$，轨迹为圆。一般情形为椭圆。取$a=b,\delta=\pi/2$时，在$t=0$电场沿$+x$，随后转向$+y$。从$+z$端朝光源看、且把$x$向右$y$向上画在纸上，此旋转为逆时针；换观察方向会改变名称，故以下均以这个时序说明旋向。

\Diagram{fig:o4states}{同一观察约定下的线、圆和椭圆轨迹}{\begin{scope}[xshift=-3cm]\draw[axis](-1.1,0)--(1.2,0);\draw[axis](0,-1.1)--(0,1.2);\draw[very thick](-.8,-.8)--(.8,.8);\node[below]at(0,-1.3){线};\end{scope}\draw[axis](-1.1,0)--(1.2,0);\draw[axis](0,-1.1)--(0,1.2);\draw[curve](0,0)circle(.9);\draw[->](.9,0)arc(0:50:.9);\node[below]at(0,-1.3){圆};\begin{scope}[xshift=3cm]\draw[axis](-1.1,0)--(1.2,0);\draw[axis](0,-1.1)--(0,1.2);\draw[curve,rotate=25](0,0)ellipse(1 and .5);\node[below]at(0,-1.3){椭圆};\end{scope}}

\subsection{消元方程与主轴}\label{th:o4ellipse}
由$E_x/a=\cos\omega t$和$E_y/b=\cos\omega t\cos\delta+\sin\omega t\sin\delta$消去时间，得
\begin{equation}\label{eq:o4ellipse}
 \frac{E_x^2}{a^2}+\frac{E_y^2}{b^2}
 -\frac{2E_xE_y\cos\delta}{ab}=\sin^2\delta.
\end{equation}
若$a$或$b$为零，应直接回到参数式，避免除零。设主轴相对$x$转角为$\psi$，令$E_\xi=E_x\cos\psi+E_y\sin\psi$、$E_\eta=-E_x\sin\psi+E_y\cos\psi$。对一个周期取平均，$\langle E_x^2\rangle=a^2/2$、$\langle E_y^2\rangle=b^2/2$、$\langle E_xE_y\rangle=ab\cos\delta/2$，于是
\[
 2\langle E_\xi E_\eta\rangle
 =ab\cos\delta\cos2\psi+\frac{b^2-a^2}{2}\sin2\psi.
\]
非退化椭圆的两主轴分量相差$\pm\pi/2$，其相关平均为零；令右边为零得到
\[
 \tan2\psi=\frac{2ab\cos\delta}{a^2-b^2}.
\]
实际应由$\operatorname{atan2}(2ab\cos\delta,a^2-b^2)$确定$2\psi$。半轴平方是矩阵$2\langle\vct E\vct E^{\mathsf T}\rangle$的两个本征值$\Lambda$，其方程为
\[
 (a^2-\Lambda)(b^2-\Lambda)-a^2b^2\cos^2\delta=0,
 \qquad \Lambda^2-(a^2+b^2)\Lambda+a^2b^2\sin^2\delta=0.
\]
解这个二次方程，长短半轴平方分别为
\[
 A_\pm^2=\frac{a^2+b^2\pm\sqrt{(a^2-b^2)^2+4a^2b^2\cos^2\delta}}2.
\]
这是时间参数轨迹的最大、最小场强平方，不是分别把$x,y$分量振幅当成半轴。只有$\cos\delta=0$或特殊对称情形，坐标轴才直接就是椭圆主轴。

\begin{example}\label{ex:o4ellipse}
Jones矢量为$(2,\ee^{i\pi/3})^{\mathsf T}E_0$。求椭圆方程、主轴方向、半轴及旋转时序。
\Diagram{fig:o4ellipseq}{两正交分量的振幅和相位差}{\draw[axis](-2.5,0)--(2.7,0)node[right]{$E_x/E_0$};\draw[axis](0,-1.5)--(0,1.7)node[above]{$E_y/E_0$};\draw[guide](-2,-1)rectangle(2,1);\node[above]at(1.5,2.3){$a=2,\ b=1,\ \delta=\pi/3$};}
\end{example}
\begin{solution}
采用实场$\operatorname{Re}(\vct E\ee^{-i\omega t})$。图\ref{fig:o4ellipseq}中$X=E_x/E_0=2\cos\omega t$、$Y=E_y/E_0=\cos(\omega t-\pi/3)$。由$Y=X/4+(\sqrt3/2)\sin\omega t$消去时间，
\[
 (Y-X/4)^2=\tfrac34(1-X^2/4)
 \quad\Longrightarrow\quad X^2/4+Y^2-XY/2=3/4.
\]
二次型矩阵为$Q=\bigl(\begin{smallmatrix}1/4&-1/4\\-1/4&1\end{smallmatrix}\bigr)$，其特征方程$\lambda^2-(5/4)\lambda+3/16=0$，得$\lambda_\pm=(5\pm\sqrt{13})/8$。
\Diagram{fig:o4ellipsea}{倾斜椭圆的长短轴；箭头表示由参数方程确定的运动方向}{\draw[axis](-2.5,0)--(2.7,0)node[right]{$X$};\draw[axis](0,-1.6)--(0,1.8)node[above]{$Y$};\draw[curve,domain=0:360,samples=100]plot({2*cos(\x)},{cos(\x-60)});\draw[guide](-2.1,-.636)--(2.1,.636);\draw[guide](-.4,1.32)--(.4,-1.32);\draw[->,thick](2,.5)--(1.97,.643);}
图\ref{fig:o4ellipsea}的半轴平方为右侧常数$3/4$除以相应特征值，故长短半轴分别满足
\[
 A^2/E_0^2=(5+\sqrt{13})/2,\qquad
 B^2/E_0^2=(5-\sqrt{13})/2.
\]
长轴对应较小的二次型特征值。旋转坐标消去交叉项给$\tan2\psi=2/3$；结合较小特征值的特征向量，取$\psi=\tfrac12\arctan(2/3)=16.85^\circ$，而非相差$90^\circ$的短轴方向。

判断旋向须回到时间参数：$t=0$处$(X,Y)=(2,1/2)$，$\dot X=0$、$\dot Y=(\sqrt3/2)\omega>0$，所以电场从右侧向上运动，与图中箭头一致。几何方程只含$\cos\delta$和$\sin^2\delta$，对$\delta$与$-\delta$给出同一几何轨迹，旋向须由时间参数确定。
\end{solution}

\section{自然光、部分偏振与检偏}\label{th:o4partial}
理想线偏振片透过轴单位矢量$\vct p$，其Jones矩阵为$P=\vct p\vct p^{\mathsf T}$。线偏振入射与轴成$\theta$时，振幅乘$\cos\theta$，光强满足马吕斯定律$I=I_0\cos^2\theta$。自然光在任意正交方向平均等分能量，故理想偏振片透过$I_0/2$，与透过轴角无关。

圆偏振经过线偏振片也透过一半，与转角无关，但出射变成线偏振。椭圆偏振转动检偏器一般有非零最大和最小值，分别正比$A_+^2,A_-^2$。自然光叠加线偏振光时，$I(\theta)=I_u/2+I_p\cos^2(\theta-\theta_0)$，所以$I_{\max}=I_u/2+I_p$、$I_{\min}=I_u/2$，偏振度$I_p/(I_u+I_p)=(I_{\max}-I_{\min})/(I_{\max}+I_{\min})$。

若光中含圆偏振分量，旋转单个线偏振片得到的只是线偏振调制度。完整偏振度还需测量圆偏振信息。定义强度单位一致的斯托克斯量$S_0=\langle|E_x|^2+|E_y|^2\rangle$、$S_1=\langle|E_x|^2-|E_y|^2\rangle$、$S_2=2\operatorname{Re}\langle E_x^*E_y\rangle$、$S_3=2\operatorname{Im}\langle E_x^*E_y\rangle$，则完全偏振$S_1^2+S_2^2+S_3^2=S_0^2$；一般偏振度为$\sqrt{S_1^2+S_2^2+S_3^2}/S_0$。

\begin{example}\label{ex:o4malus}
两理想偏振片透过轴正交，自然光先经过第一片，再在中间插入轴与第一片成$\theta$角的第三片。求透过强度及最优角度。
\Diagram{fig:o4malusq}{三片顺序为$0,\theta,\pi/2$}{\foreach\x/\t in{0/0,2/\theta,4/{\pi/2}}{\draw[curve](\x,-1)--(\x,1);\node[above]at(\x,1){$\t$};}\draw[axis](-1,0)--(5,0)node[right]{光};}
\end{example}
\begin{solution}
图\ref{fig:o4malusq}中，自然光的横向偏振方向均匀分布，第一片透过强度为$I_0\langle\cos^2\varphi\rangle=I_0/2$。它之后的光沿第一片轴线偏振，不再是自然光。中间片与该轴成$\theta$角，透过场投影乘$\cos\theta$，强度乘$\cos^2\theta$；最后片与中间轴夹角为$\pi/2-\theta$，再乘$\sin^2\theta$，故
\[
 I=\frac{I_0}{2}\cos^2\theta\sin^2\theta
 =\frac{I_0}{8}\sin^22\theta.
\]
\Diagram{fig:o4malusa}{逐次投影后偏振方向改变，不能把中间片视为单纯削弱光强}{\draw[axis](0,0)--(2.2,0)node[right]{$P_1$};\draw[axis](0,0)--(1.6,1.6)node[right]{$P_2$};\draw[axis](0,0)--(0,2.2)node[above]{$P_3$};\draw[guide](2,0)--(1,1)--(0,1);}
图\ref{fig:o4malusa}按传播次序标为$P_1,P_2,P_3$，中间插入片为$P_2$，末片为$P_3$。由$\sin^22\theta\leq1$，最大透过强度为$I_0/8$，取$\theta=\pi/4$模$\pi/2$；在$0\leq\theta\leq\pi/2$内唯一最优角为45度。

逐段能量也可核对：最优时$I_1=I_0/2$、$I_2=I_0/4$、$I_3=I_0/8$，每片都吸收部分功率。没有中间片时两轴正交，最后为零；插入后中间片把输出偏振改为45度，使它在最后轴上有非零投影。透过增加来自偏振态改变，整个过程仍有$I_3<I_2<I_1<I_0$，并未由吸收元件产生能量。
\end{solution}

\section{单轴晶体中的普通光与非普通光}\label{th:o4crystal}
各向异性介质的线性关系为$\vct D=\varepsilon_0\boldsymbol\varepsilon\vct E$，主轴系中介电张量可对角化。一般$\vct D$与$\vct E$不平行。无自由电荷仍要求$\vct k\cdot\vct D=0$，却不能直接推为$\vct k\cdot\vct E=0$。

单轴晶体取光轴为$z$，$\boldsymbol\varepsilon=\operatorname{diag}(n_o^2,n_o^2,n_e^2)$，这里$n_e$是传播方向垂直光轴时的非普通主折射率。波矢与光轴夹角为$\vartheta$。普通光电场垂直由$\vct k$和光轴张成的主截面，折射率恒为$n_o$。非普通光的$\vct D$在主截面内并垂直$\vct k$，其色散曲面为
\[
 \frac{k_x^2}{n_e^2}+\frac{k_z^2}{n_o^2}=\frac{\omega^2}{c^2},
\]
所以有效折射率
\[
 \frac1{n_{e,\rm eff}^2(\vartheta)}=
 \frac{\cos^2\vartheta}{n_o^2}+\frac{\sin^2\vartheta}{n_e^2}.
\]
沿光轴两种波简并，都为$n_o$；垂直光轴时分别为$n_o,n_e$。非普通光的能流方向由群速度或坡印廷矢量决定，一般不平行于波矢，产生走离。相位匹配用波矢，光线方向与光束分离则要用能流，二者不能混写。

\Diagram{fig:o4crystal}{波矢与位移矢量垂直，电场和能流可偏斜}{\draw[axis](-.3,0)--(3.4,0)node[right]{$x$};\draw[axis](0,-.3)--(0,3.3)node[above]{光轴$z$};\draw[axis](0,0)--(2,2.4)node[right]{$\vct k$};\draw[axis](0,0)--(-1.8,1.5)node[left]{$\vct D_e$};\draw[axis](0,0)--(-1,1.9)node[above]{$\vct E_e$};\draw[guide](0,0)--(2.85,1.5)node[right]{$\vct S_e$};}

\section{波片及偏振态变换}\label{th:o4plates}
波片正常入射，两个本征偏振轴在片面内，取快轴折射率$n_f$、慢轴$n_s>n_f$，厚度$h$。删去共同相位$\exp(ik_0n_fh)$后，Jones矩阵为$J=\operatorname{diag}(1,\ee^{i\delta})$，$\delta=2\pi(n_s-n_f)h/\lambda_0$。四分之一波片$\delta=\pi/2$模$2\pi$；半波片$\delta=\pi$模$2\pi$。共同相位不影响单束偏振，但在干涉仪中仍可能影响两臂相位。

如果波片快轴相对实验室$x$转角$\alpha$，先用$R(-\alpha)$把分量投到快慢轴，再乘对角矩阵，最后转回：$J(\alpha)=R(\alpha)\operatorname{diag}(1,\ee^{i\delta})R(-\alpha)$，其中$R(\alpha)=\bigl(\begin{smallmatrix}\cos\alpha&-\sin\alpha\\\sin\alpha&\cos\alpha\end{smallmatrix}\bigr)$。矩阵顺序代表实际操作，交换两片的顺序一般改变结果。

\begin{example}\label{ex:o4quarter}
线偏振入射与四分之一波片快轴成$\theta$，求输出Jones矢量、椭圆轴及旋转方向，特别讨论$\theta=45^\circ$。
\Diagram{fig:o4quarterq}{输入线偏振在快慢轴上的投影}{\draw[axis](0,0)--(3,0)node[right]{快$x$};\draw[axis](0,0)--(0,2.8)node[above]{慢$y$};\draw[axis](0,0)--(2,2)node[right]{输入};\node at(.9,.35){$\theta$};}
\end{example}
\begin{solution}
图\ref{fig:o4quarterq}输入为$E_0(\cos\theta,\sin\theta)^{\mathsf T}$，输出为$E_0(\cos\theta,i\sin\theta)^{\mathsf T}$。两个分量正交相差$\pi/2$，所以主轴就是快慢轴，半轴绝对值$E_0|\cos\theta|$和$E_0|\sin\theta|$。
\Diagram{fig:o4quartera}{45度输入产生圆偏振，固定位置由$+x$转向$+y$}{\draw[axis](-1.5,0)--(1.7,0)node[right]{$x$};\draw[axis](0,-1.5)--(0,1.7)node[above]{$y$};\draw[curve](0,0)circle(1.1);\draw[->,thick](1.1,0)arc(0:65:1.1);\node[right]at(2,.5){$(1,i)/\sqrt2$};}
若$0<\theta<\pi/2$，固定位置$t=0$电场沿$+x$，随时间向$+y$转，图\ref{fig:o4quartera}按此方向画箭头。$\theta=45^\circ$为圆；$\theta=0,90^\circ$仍是线偏振。把片转90度等效交换快慢轴，会反转圆偏振旋向，不能只说“经过四分之一波片必变圆”。
\end{solution}

\begin{example}\label{ex:o4half}
入射线偏振方向相对实验室$x$为$\theta$，半波片快轴方向为$\alpha$。求出射偏振方向，并求将偏振方向由10度转到70度的快轴位置。
\Diagram{fig:o4halfq}{半波片快轴与输入偏振的夹角}{\draw[axis](0,0)--(3,0)node[right]{$x$};\draw[axis](0,0)--(2.5,1)node[right]{$\theta$};\draw[guide](0,0)--(1.6,2.5)node[above]{$\alpha$};}
\end{example}
\begin{solution}
在图\ref{fig:o4halfq}的快慢轴基底，输入为$(\cos(\theta-\alpha),\sin(\theta-\alpha))$。半波片把慢轴分量反号，成为$(\cos(\theta-\alpha),-\sin(\theta-\alpha))$，这仍是线偏振，方向在快轴另一侧。
\Diagram{fig:o4halfa}{半波片把输入方向关于快轴作镜面对称}{\draw[axis](0,0)--(3,0)node[right]{$x$};\draw[axis](0,0)--(2.8,.49)node[right]{10度};\draw[guide](0,0)--(2.3,1.93)node[right]{快轴40度};\draw[axis](0,0)--(1.03,2.82)node[above]{70度};}
转回实验室，输出方向$\theta_o=2\alpha-\theta$模$\pi$。故图\ref{fig:o4halfa}可取$\alpha=40^\circ$；快轴再转90度得到同一线偏振方向但共同相位可能不同。半波片旋转$\Delta\alpha$，输出线偏振方向转$2\Delta\alpha$。
\end{solution}

\begin{example}\label{ex:o4identify}
未知光可能为完全线、圆、椭圆偏振，或自然光。给出理想线偏振片和可转四分之一波片的鉴别程序，说明部分偏振光需要怎样补充。
\Diagram{fig:o4identifyq}{先测线检偏响应，再加入可转波片}{\draw[axis](-1,0)--(6,0);\draw[curve](0,-1)--(0,1);\draw[curve](2,-1)--(2,1);\draw[curve](4,-1)--(4,1);\node[above]at(0,1){波片};\node[above]at(2,1){检偏器};\node[above]at(4,1){功率计};}
\end{example}
\begin{solution}
先移去图\ref{fig:o4identifyq}的波片，转检偏器：周期变化且能消光为完全线偏振；变化但不能消光可能为完全椭圆偏振，也可能为部分偏振；不变化可能是圆偏振或自然光，不能立即下结论。
\Diagram{fig:o4identifya}{波片把相位信息转成可由线检偏器测出的强度差}{\draw[curve](0,0)rectangle(2.2,1.2);\node at(1.1,.6){相位差};\draw[->](2.3,.6)--(3.3,.6);\draw[curve](3.5,0)rectangle(6,1.2);\node at(4.75,.6){旋转后的强度差};}
加入四分之一波片，转片和检偏器：圆偏振可变成线偏振而消光，自然光经过理想无损波片仍不随检偏角变化。完全椭圆偏振可把波片轴对准椭圆主轴，抵消$\pm\pi/2$相差而变线。图\ref{fig:o4identifya}说明波片提供的是原来缺少的相位信息。一般部分偏振不能靠“是否消光”完整确定，应测$0,90,45,135$度线偏振强度，以及两种圆偏振分析强度，构成$S_0,S_1,S_2,S_3$。若入射光含非零自然成分，仅转动这块无损波片和随后的一片检偏器，不能使总透射光完全消失。
\end{solution}

\section{界面边界条件与菲涅耳公式}\label{th:o4fresnel}
取界面$z=0$，介质1在$z<0$、介质2在$z>0$，无自由面电荷和面电流。入射面为$xz$面，入射、反射、透射波矢单位方向依次为$(\sin i,0,\cos i)$、$(\sin i,0,-\cos i)$、$(\sin r,0,\cos r)$。沿界面相位匹配给$n_1\sin i=n_2\sin r$，反射角等于入射角。

以下取无损非磁性介质。$s$偏振统一取电场单位矢量$\hat{\vct y}$；每条光线的$p$单位矢量取$\hat{\vct s}\times\hat{\vct k}$，使$\hat{\vct k}\times\hat{\vct p}=\hat{\vct s}$。于是入射$p$电场为$(\cos i,0,-\sin i)$，反射为$(-\cos i,0,-\sin i)$，两者在正入射极限的切向基方向相反。

\subsection{$s$偏振}
切向电场连续给$1+r_s=t_s$。磁场$\vct H=(n/\mu_0c)\hat{\vct k}\times\vct E$，切向磁场连续给$n_1\cos i(1-r_s)=n_2\cos r\,t_s$。消元得
\[
 r_s=\frac{n_1\cos i-n_2\cos r}{n_1\cos i+n_2\cos r},\quad
 t_s=\frac{2n_1\cos i}{n_1\cos i+n_2\cos r}.
\]
\subsection{$p$偏振}
在上述局部基约定下，切向电场给$\cos i(1-r_p)=\cos r\,t_p$，磁场给$n_1(1+r_p)=n_2t_p$。由后式取$t_p=n_1(1+r_p)/n_2$代入前式，整理为
\[
 n_2\cos i-n_1\cos r=(n_2\cos i+n_1\cos r)r_p,
\]
再回代求$t_p$，得
\[
 r_p=\frac{n_2\cos i-n_1\cos r}{n_2\cos i+n_1\cos r},\quad
 t_p=\frac{2n_1\cos i}{n_2\cos i+n_1\cos r}.
\]
正入射时$r_s=-r_p$，只是反射$p$基反向的结果；用同一固定切向电场比较，两物理反射系数相同。若$\mu_r\ne1$，磁场系数$n$应改用相应波导纳，不能把非磁性公式直接照搬。

\section{能流、互易与布儒斯特角}\label{th:o4energy}
反射率是反射法向能流与入射法向能流之比，故$R=|r|^2$。透射率为
\[
 T=\frac{n_2\cos r}{n_1\cos i}|t|^2.
\]
无损界面在存在传播透射波时$R+T=1$。若只取沿各自光线定义的强度比$I_t/I_i=(n_2/n_1)|t|^2$，斜入射时它与$R$之和一般不为1，因为比较的截面面积不同，不能据此宣称斜入射能量不守恒。

\begin{example}\label{ex:o4stokes}
对$s$偏振，正向界面振幅系数为$r,t$，反向入射为$r',t'$。证明$r'=-r$、$tt'=1-r^2$，并解释斯托克斯关系中各束光的相位方向。
\Diagram{fig:o4stokesq}{同一界面的正反向散射}{\draw[curve](-2,0)--(2,0);\draw[->](-1.5,-1.5)--(0,0);\draw[->](0,0)--(1.5,-1.5)node[right]{$r$};\draw[->](0,0)--(1.1,1.5)node[right]{$t$};\node[left]at(-2,-.6){$n_1$};\node[left]at(-2,.6){$n_2$};}
\end{example}
\begin{solution}
图\ref{fig:o4stokesq}取$s$偏振的共同切向电场方向作为振幅基底，令$a=n_1\cos i$、$b=n_2\cos r$。电场、磁场切向连续分别给$1+r=t$、$a(1-r)=bt$。将首式代入次式，解得
\[
 r=\frac{a-b}{a+b},\qquad t=\frac{2a}{a+b}.
\]
反向入射时交换$a,b$，因此$r'=(b-a)/(a+b)=-r$、$t'=2b/(a+b)$。
\Diagram{fig:o4stokesa}{时间反演的两束入射共同恢复原来的入射光}{\draw[curve](-2,0)--(2,0);\draw[->](-1.5,-1.5)--(0,0)node[midway,left]{$r$};\draw[->](1.1,1.5)--(0,0)node[midway,right]{$t$};\draw[->](0,0)--(1.5,-1.5)node[midway,right]{$1$};\node[above]at(-1.6,1.5){另一输出相消};}
图\ref{fig:o4stokesa}的两束反向入射按原反、透射振幅$r,t$取值。在恢复原入射光的输出端，合振幅为$r\,r+t\,t'$；另一端为$r\,t+t\,r'$。代入上述系数，
\[
 r^2+tt'=\frac{(a-b)^2+4ab}{(a+b)^2}=1,
 \qquad rt+tr'=t(r+r')=0.
\]
所以原波被恢复，另一输出相消。$r'=-r$表示反向反射相位相差$\pi$，不是反射率变成负数。能量透射率应为$T=(b/a)|t|^2$，由此$R+T=1$；不能直接令$T=t^2$。这里的时间反演图解采用无损传播及实系数的两侧入射端口。全反射的倏逝支、复系数或吸收介质须连同共轭、通道与能流条件重新处理。
\end{solution}

\subsection{布儒斯特角}\label{th:o4brewster}
$p$反射为零需$n_2\cos i=n_1\cos r$，再与$n_1\sin i=n_2\sin r$联立，非平凡折射率差时得到$i+r=\pi/2$，所以$\tan i_B=n_2/n_1$。这时反射光只含$s$偏振。若$n_1=n_2$，所有角度都没有界面反射，不能把它误解成只有45度一个布儒斯特角。

地面或水面的近布儒斯特反射主要为平行表面的$s$偏振，所以让眼镜透过轴接近竖直可削弱这部分眩光。实际粗糙表面、金属和吸收介质不能全部用理想介质界面的单一布儒斯特角描述。

\section{全反射相位与倏逝波}\label{th:o4tir}
设$n_1>n_2$，入射角超过$i_c=\arcsin(n_2/n_1)$。切向波矢$k_x=k_0n_1\sin i$不变，介质2的法向分量满足$k_{z2}^2=k_0^2n_2^2-k_x^2<0$。介质2位于$z>0$，为使场远离界面不发散，取$k_{z2}=i\kappa$，其中
\[
 \kappa=k_0\sqrt{n_1^2\sin^2i-n_2^2},\quad
 E_t\propto\ee^{ik_xx-\kappa z-i\omega t}.
\]
振幅衰减长度为$1/\kappa$，强度衰减长度为$1/(2\kappa)$。令$q=\sqrt{n_1^2\sin^2i-n_2^2}$，在菲涅耳系数中取$n_2\cos r=iq$，得
\[
 r_s=\frac{n_1\cos i-iq}{n_1\cos i+iq},\quad
 r_p=\frac{n_2^2\cos i-in_1q}{n_2^2\cos i+in_1q}.
\]
分子分母互为共轭，模均为1，但相位依角度和偏振而不同。全反射可改变偏振态；“反射率1”不等于“相位不变”。单界面介质2中存在倏逝场，时间平均法向能流为零，但沿界面可有非零能流。

\begin{example}\label{ex:o4tir}
玻璃$n_1=1.5$向空气$n_2=1$以45度入射，波长$\lambda_0$。求$s,p$反射相位及倏逝振幅衰减长度，并判断能否把一次反射等同于四分之一波片。
\Diagram{fig:o4tirq}{超过临界角后无传播折射光，但有衰减场}{\draw[curve](-2.5,0)--(2.5,0);\draw[guide](0,-2)--(0,2);\draw[->](-1.8,-1.8)--(0,0);\draw[->](0,0)--(1.8,-1.8);\foreach\y/\l in{.3/2,.7/1.3,1.1/.8,1.5/.5}{\draw[curve](-\l,\y)--(\l,\y);}\node[right]at(2,.7){倏逝场};}
\end{example}
\begin{solution}
图\ref{fig:o4tirq}中$q=\sqrt{2.25/2-1}=1/(2\sqrt2)$，$n_1\cos i=3/(2\sqrt2)$，所以$\phi_s=-2\arctan(1/3)\simeq-36.87^\circ$。$p$系数相位为$\phi_p=-2\arctan[n_1q/(n_2^2\cos i)]=-2\arctan(3/4)\simeq-73.74^\circ$。
\Diagram{fig:o4tira}{两偏振的单位模反射系数具有不同相位}{\draw[axis](-.3,0)--(2.6,0)node[right]{实部};\draw[axis](0,-2.5)--(0,.5)node[above]{虚部};\draw[curve](2,0)arc(0:-90:2);\draw[->,thick](0,0)--(1.6,-1.2)node[right]{$r_s$};\draw[->,thick](0,0)--(.56,-1.92)node[below]{$r_p$};}
图\ref{fig:o4tira}的相对相差为$-36.87^\circ$，不是$-90^\circ$，所以一次反射不等于四分之一波片。衰减常数$\kappa=2\pi q/\lambda_0=\pi/(\sqrt2\lambda_0)$，振幅长度$\ell=\sqrt2\lambda_0/\pi$。在给定基向量约定下比较两分量相位，还需同步追踪反射$p$基与实验室固定横向轴的关系。
\end{solution}

\begin{example}\label{ex:o4frustrated}
两块相同高折射率玻璃间隔空气层厚度$d$，入射角在单个玻璃—空气界面满足全反射。说明为何第二块玻璃足够靠近时可有透射，并给出厚间隙极限下的主要衰减尺度。
\Diagram{fig:o4frustratedq}{有限低折射率层中的倏逝耦合；间隙内不画传播光线}{
\draw[curve](-2.2,0)--(2.4,0);\draw[curve](-2.2,1.1)--(2.4,1.1);
\node[left]at(-2.2,-.6){$n_1$};\node[left]at(-2.2,.55){$n_2$};\node[left]at(-2.2,1.7){$n_1$};
\draw[->](-1.6,-1.6)--(0,0);\draw[->](0,0)--(1.6,-1.6);
\draw[->](.8,1.1)--(2.1,2.4);
\foreach\y/\l in{.23/1.1,.51/.8,.82/.55}{\draw[densely dotted](-\l,\y)--(\l,\y);}
\node[right]at(1.1,.55){倏逝场};
\draw[<->](2.8,0)--(2.8,1.1)node[midway,right]{$d$};
}
\end{example}
\begin{solution}
图\ref{fig:o4frustratedq}间隙内仍有相同实切向波矢和虚法向波矢$\pm i\kappa$。单界面半无限介质只许衰减支；有限间隙必须同时满足两端边界，场通常含$A\ee^{-\kappa z}+B\ee^{\kappa z}$。增长支在有限区间不发散，不能提前删去。
\Diagram{fig:o4frustrateda}{有限层两指数支叠加，可把能流耦合到右侧传播波}{\draw[axis](0,0)--(5.4,0)node[right]{$z$};\draw[guide](1,0)--(1,2.8);\draw[guide](4,0)--(4,2.8);\draw[curve,domain=1:4,samples=60]plot(\x,{2.2*exp(-.65*(\x-1))});\node[below]at(1,0){$0$};\node[below]at(4,0){$d$};\draw[->](4.2,.6)--(5,.6);}
先对$s$偏振完成匹配。令$k_0=2\pi/\lambda_0$、$a=n_1\cos i$、$b=\sqrt{n_1^2\sin^2i-n_2^2}$，于是$\kappa=k_0b$。把右界面场幅记为$t$，右侧只有出射波，故$E(d)=t$、$E'(d)=ik_0at$。间隙方程$E''=\kappa^2E$给
\[
 E(z)=t\left[\cosh\kappa(d-z)
 -i\frac ab\sinh\kappa(d-z)\right].
\]
它满足右界面的函数、导数两个条件。左侧入射幅归一为1，有$E(0)=1+r$、$E'(0)=ik_0a(1-r)$。代入上式，把导数条件除以$ik_0a$，再与场条件相加，得到
\[
 2=t\left[2\cosh\kappa d
 +i\left(\frac ba-\frac ab\right)\sinh\kappa d\right].
\]
左右传播介质相同，透射率为$|t|^2$。使用$\cosh^2u=1+\sinh^2u$，整理为
\[
 T_s=\left[1+\frac{(a^2+b^2)^2}{4a^2b^2}
 \sinh^2\kappa d\right]^{-1}.
\]
对$p$偏振，以磁场幅为未知量，界面连续的是$H$和$H'/n^2$。相同消元给上式的$a=\cos i/n_1$、$b=\sqrt{n_1^2\sin^2i-n_2^2}/n_2^2$，指数中的$\kappa$仍取原值。

当$d\to0$，两块同种玻璃接合，$T\to1$。当$\kappa d\gg1$，$\sinh\kappa d\simeq\ee^{\kappa d}/2$，所以
\[
 T\simeq\frac{16a^2b^2}{(a^2+b^2)^2}\ee^{-2\kappa d}.
\]
图\ref{fig:o4frustrateda}示意厚间隙中的主要衰减包络。单界面零法向平均能流的结论针对单一衰减支；有限层中两支场的交叉项可以传递实的法向平均能流，第二块玻璃因而接收到传播波。
\end{solution}
\section{偏振椭圆的代数结构}
\subsection{相位差怎样进入实轨迹}
令$X=E_x/a$、$Y=E_y/b$，则$X=\cos\omega t$、$Y=X\cos\delta+\sin\omega t\sin\delta$。将$Y-X\cos\delta$平方，等于$(1-X^2)\sin^2\delta$，移项得到$X^2+Y^2-2XY\cos\delta=\sin^2\delta$。这一步保留了轨迹几何，却把$\sin\delta$的正负平方掉，所以旋向必须另从时间参数恢复。

场强平方$E_x^2+E_y^2$可用二倍角展开为
\[
 \frac{a^2+b^2}2+
 \frac{a^2+b^2\cos2\delta}2\cos2\omega t
 +\frac{b^2\sin2\delta}2\sin2\omega t.
\]
后两项组成单一正弦，其振幅是$\tfrac12\sqrt{a^4+b^4+2a^2b^2\cos2\delta}$。因此最大、最小场强平方为$(a^2+b^2\pm\sqrt{a^4+b^4+2a^2b^2\cos2\delta})/2$，与前述主轴半径平方一致。椭圆长短半轴因而可由场模随时间变化的极大、极小值确定。

\subsection{退化情形的处理}
$\delta=0,\pi$时短半轴为零，轨迹退为直线；若$a=b$且$\delta=\pm\pi/2$，长短半轴相等，轨迹为圆，主轴方向任意。圆偏振时$\tan2\psi$的表达式成为$0/0$，正是因为圆上不存在唯一的主轴方向。

圆偏振的Jones矢量乘任意共同相位，固定位置上电场的起始角会改变，但旋转方向不变。同一椭圆旋向也不由相位$\delta$的主值名称决定，需同时说明时间指数约定、光传播方向和观察方向。所有图中的箭头均按本书$\ee^{-i\omega t}$与从$+z$端观察的约定作出。

\section{从Jones矢量到相干矩阵}\label{th:o4coherence}
对一般光束定义$2\times2$相干矩阵$C=\langle\vct e\vct e^\dagger\rangle$。它的对角元是两个正交分量的平均强度，非对角元保留它们的相关相位。对任何分析矢量$\vct a$，透射强度为$\vct a^\dagger C\vct a\geq0$，所以$C$是半正定厄米矩阵。

用本书的斯托克斯约定，
\[
 C=\begin{pmatrix}\langle|E_x|^2\rangle&\langle E_xE_y^*\rangle\\
 \langle E_yE_x^*\rangle&\langle|E_y|^2\rangle\end{pmatrix}
 =\frac12\begin{pmatrix}
 S_0+S_1&S_2-iS_3\\S_2+iS_3&S_0-S_1
 \end{pmatrix}.
\]
记$s^2=S_1^2+S_2^2+S_3^2$。将本征方程的行列式展开，
\begin{align*}
 0=\det(C-\lambda I)
 &=\left(\frac{S_0+S_1}{2}-\lambda\right)
 \left(\frac{S_0-S_1}{2}-\lambda\right)-\frac{S_2^2+S_3^2}{4}\\
 &=\lambda^2-S_0\lambda+\frac{S_0^2-s^2}{4}.
\end{align*}
二次方程给$\lambda_\pm=(S_0\pm s)/2$。半正定等价于$S_0\geq s\geq0$，故$S_1^2+S_2^2+S_3^2\leq S_0^2$。对非零光束定义$P=s/S_0$，则
\[
 P=\frac{\lambda_+-\lambda_-}{\lambda_++\lambda_-},\qquad
 P^2=1-\frac{4\det C}{(\operatorname{tr}C)^2}.
\]
完全偏振时$\lambda_-=0$、矩阵秩为1；自然光$C=(S_0/2)I$，两个本征值相同。对一般部分偏振，还可写$C=\lambda_-I+(\lambda_+-\lambda_-)\vct u_+\vct u_+^\dagger$：第一项是自然光，第二项是沿本征态$\vct u_+$的完全偏振分量。

\Diagram{fig:o4coherencematrix}{纯偏振是秩1相干矩阵，自然光平均后没有固定交叉相位}{\draw[curve](0,0)rectangle(2.6,1.2);\node at(1.3,.6){$C=\vct e\vct e^\dagger$};\draw[curve](3.7,0)rectangle(6.3,1.2);\node at(5,.6){$C=(S_0/2)I$};\node[above]at(1.3,1.3){完全偏振};\node[above]at(5,1.3){自然光};}
图\ref{fig:o4coherencematrix}也解释为何自然光不能选一个Jones矢量表示。若把两个正交确定场以固定相位直接相加，就得到某一种完全偏振态，而非自然光；自然光需要对随机相位或随机方向作统计平均。

无损波片的Jones矩阵$J$为酉矩阵，输出相干矩阵$C'=JCJ^\dagger$。本征值与迹不变，所以总强度和完整偏振度不变。波片可以在斯托克斯空间旋转偏振方向和椭圆率，却不能凭空把自然光变成完全偏振光。偏振片为投影矩阵，会丢弃一部分能量，才可以提高透射束偏振度。

\section{单轴晶体色散式的推导}
取$\vct k=(k_x,0,k_z)$、介电张量$\boldsymbol\varepsilon=\operatorname{diag}(n_o^2,n_o^2,n_e^2)$，假定非磁性。麦克斯韦方程给
\[
 \vct k\times(\vct k\times\vct E)+k_0^2\boldsymbol\varepsilon\vct E=0.
\]
若$\vct E$沿$y$，得到$k_x^2+k_z^2=k_0^2n_o^2$，即普通光。若$\vct E$在$xz$面，则
\[
 \begin{pmatrix}k_0^2n_o^2-k_z^2&k_xk_z\\
 k_xk_z&k_0^2n_e^2-k_x^2\end{pmatrix}
 \begin{pmatrix}E_x\\E_z\end{pmatrix}=0.
\]
要有非零场，行列式为零。展开、消去$k_x^2k_z^2$项，再除以$k_0^2n_o^2n_e^2$，得到$k_x^2/n_e^2+k_z^2/n_o^2=k_0^2$，就是非普通色散椭圆。

高斯定律给$k_xn_o^2E_x+k_zn_e^2E_z=0$，所以$\vct D\perp\vct k$，但$\vct E$不一定垂直。无损非色散介质的群速度方向为等频曲面的法线，通常也就是能流方向；非普通光因此发生走离。沿光轴$k_x=0$时，两支都给折射率$n_o$，横向任意线偏振可作为等价基底，不会发生通常意义的双折射分离。

\section{菲涅耳能量关系的逐项证明}
对$s$偏振，令$a=n_1\cos i$、$b=n_2\cos r$，$r_s=(a-b)/(a+b)$、$t_s=2a/(a+b)$。法向透射率$T_s=(b/a)t_s^2=4ab/(a+b)^2$，故$R_s+T_s=[(a-b)^2+4ab]/(a+b)^2=1$。

对$p$偏振，令分母$D=n_2\cos i+n_1\cos r$，则
\[
 1-r_p^2=\frac{4n_1n_2\cos i\cos r}{D^2},
\]
而法向透射率
\[
 T_p=\frac{n_2\cos r}{n_1\cos i}
 \frac{4n_1^2\cos^2i}{D^2}
 =\frac{4n_1n_2\cos i\cos r}{D^2}.
\]
因此两偏振都严格满足$R+T=1$，无需限定正入射。若取沿透射方向的强度比$I_t/I_i$而未乘$\cos r/\cos i$，相当于比较两种不同截面，当然不能直接与$R$相加作界面能量收支。

\Diagram{fig:o4fluxarea}{同一界面面积$A$对入射和透射光束的投影面积分别为$A\cos i,A\cos r$}{\draw[curve](-2.5,0)--(2.5,0);\draw[->](-1.7,-1.7)--(0,0)--(1,1.8);\draw[guide](0,-2)--(0,2.3);\node[left]at(-1,-1.2){$I_i\cos i$};\node[right]at(.8,1.1){$I_t\cos r$};\node[below]at(1.8,-.2){界面$A$};}
图\ref{fig:o4fluxarea}以固定界面面积作为收支基准。对于吸收介质，透射角和波阻抗可能变复，能流方向及损耗需要用复坡印廷矢量重新求出；本讲的实角无损公式不能直接把$n$换成复数后仍以同一实余弦解释。

\section{全反射的解析支与相位连续性}
当入射角从临界角以下增至以上，透射法向波矢由正实变为正虚。正虚支使$z>0$侧的$\ee^{ik_{z2}z}$衰减；负虚支则使场在无限远增长，与无穷远无输入的边界条件不符。若低折射率介质改在另一半空间，应重新按远离界面的衰减条件选取虚部符号。

反射系数在临界角以上沿复平面的单位圆变化。其模恒1，能量全部反回，但$s,p$相位差不恒定。若入射为两分量都有的线偏振，一次全反射一般变椭圆；两次适当角度的反射可以累积为四分之一波延迟。多次反射还要追踪各入射面如何转动，局部$s,p$基底之间若有几何旋转，就不能只把标量相位差相加。

倏逝场中$\vct E,\vct H$某些分量相差90度，平均法向坡印廷流为零，同时有非零瞬时能量往返交换。有限空气隙含两支指数解，交叉项能产生实的法向平均流，因此受抑全反射不违反单界面结果；它改变的是边界问题本身。有限间隙的透射率应由两端边界共同求出。

\subsection{部分偏振的相关矩阵与可实现条件}
第\ref{th:o4coherence}节已由相干矩阵的本征方程证明$S_1^2+S_2^2+S_3^2\leq S_0^2$。这个不等式也给实验数据的可实现条件：固定非零总强度$S_0$后，归一化矢量$(S_1,S_2,S_3)/S_0$位于单位球内，距球心的长度就是$P$。球面对应完全偏振，球心对应自然光。若重建值落在球外，应回查分析器标定、背景扣除或光源漂移；它不能作为某种“超完全偏振”状态。$S_0=0$时无光，不定义偏振度。
\Diagram{fig:o4stokesball}{固定总强度下的归一化斯托克斯矢量；球面对应完全偏振，球内对应部分偏振}{\draw[axis](-2,0)--(2.4,0)node[right]{$S_1/S_0$};\draw[axis](0,-1.8)--(0,2)node[above]{$S_3/S_0$};\draw[curve](0,0)circle(1.5);\draw[guide](0,0)ellipse(1.5 and .38);\draw[->,thick](0,0)--(.8,.7)node[right]{$P<1$};\fill(-.9,1.2)circle(1.5pt)node[left]{$P=1$};\node[below]at(.2,-1.9){$S_2$轴的投影略去};}
图\ref{fig:o4stokesball}把总强度固定后，偏振态表示为球内或球面点。只有线检偏器时能测得$S_0,S_1,S_2$，无法确定$S_3$；必须加入波片形成圆分析器，才能恢复全部相关信息。线检偏器扫描的可见度不等于总体偏振度；圆偏振的线分析强度与方向无关，但$P=1$。

对多束互不相干光，交叉相关项平均为零，故$C_{\rm tot}=\sum_jC_j$，四个斯托克斯量也分别相加。偏振度含矢量和的模，应在相加后计算：相反偏振方向可以抵消，即使每束的$P$都为1，总光也可能为自然光。无损波片与理想检偏器对$C$的作用已在第\ref{th:o4coherence}节区分，不能用偏振度的算术平均代替这一矩阵运算。

\section{从测量强度还原偏振态}
\begin{example}\label{ex:o4stokesmeasure}
同一未知光束经六种理想分析器测得强度：水平6、竖直4、45度7、135度3、正圆偏振分析6、负圆偏振分析4。圆分析正态定义为Jones矢量$(1,i)^{\mathsf T}/\sqrt2$。求斯托克斯量、完整偏振度以及线偏振扫描的可见度。
\Diagram{fig:o4measureq}{同一光束依次采用六种偏振分析设置}{\draw[curve](0,0)rectangle(2.2,1);\node at(1.1,.5){未知光};\draw[->](2.4,.5)--(3.3,.5);\node[right]at(3.5,.5){$H,V,D,A,+,-$};}
\end{example}
\begin{solution}
图\ref{fig:o4measureq}三组互补强度之和都为10，彼此相容。故$S_0=10$，$S_1=I_H-I_V=2$，$S_2=I_D-I_A=4$，$S_3=I_+-I_-=2$。完整偏振度$\mathcal P=\sqrt{2^2+4^2+2^2}/10=\sqrt6/5\simeq.4899$。
\Diagram{fig:o4measurea}{完整偏振矢量的长度大于其线偏振平面投影}{\draw[axis](0,0)--(3.4,0)node[right]{$S_1$};\draw[axis](0,0)--(0,3)node[above]{$S_2$};\draw[axis](0,0)--(-1.5,-1.2)node[left]{$S_3$};\draw[->,thick](0,0)--(1.8,2.2);\draw[guide](1.8,0)--(1.8,2.2);\node[right]at(2.1,1){三维分量示意};}
旋转线检偏器强度$I(\theta)=\tfrac12[S_0+S_1\cos2\theta+S_2\sin2\theta]$，调制振幅$\sqrt{20}/2$，所以可见度$\sqrt{20}/10=1/\sqrt5\simeq.4472$。图\ref{fig:o4measurea}说明单线检偏只看$S_1,S_2$投影，忽略$S_3$会低估完整偏振度。此组数据含不可定相的自然成分，不能指定一个唯一Jones矢量代表全部光。
\end{solution}

\begin{example}\label{ex:o4quarteranalysis}
证明快轴45度、慢轴135度的理想四分之一波片接水平检偏器，可测得正圆分析强度$S_0/2+S_3/2$。说明转片或交换快慢轴怎样测另一旋向。
\Diagram{fig:o4analysisq}{四分之一波片把圆基分量转成正交线基分量}{\draw[axis](-1,0)--(5.5,0);\draw[curve](0,-1)--(0,1);\draw[curve](3,-1)--(3,1);\node[above]at(0,1){$Q(45^\circ)$};\node[above]at(3,1){水平检偏};}
\end{example}
\begin{solution}
图\ref{fig:o4analysisq}中，快轴沿45度、慢轴沿135度，令快轴相位为零、慢轴增相$\pi/2$，波片矩阵为
\[
 Q_{45}=\frac12\begin{pmatrix}1+i&1-i\\1-i&1+i\end{pmatrix}.
\]
后接水平检偏器只保留第一行，因此$E'_H=[(1+i)E_x+(1-i)E_y]/2=(1+i)(E_x-iE_y)/2$。公共因子模平方为$1/2$，故
\[
 I_H'=\tfrac12\langle|E_x-iE_y|^2\rangle.
\]
\Diagram{fig:o4analysisa}{两圆偏振基在该波片后分别成为水平、竖直线偏振}{\node at(0,1){$(1,i)$};\node at(0,-.3){$(1,-i)$};\draw[->](1,1)--(2.3,1);\draw[->](1,-.3)--(2.3,-.3);\node[right]at(2.5,1){水平};\node[right]at(2.5,-.3){竖直};}
按本讲$S_3=-2\operatorname{Im}\langle E_xE_y^*\rangle$的约定，展开模平方，
\begin{align*}
 \langle|E_x-iE_y|^2\rangle
 &=\langle|E_x|^2+|E_y|^2+iE_xE_y^*-iE_yE_x^*\rangle\\
 &=S_0-2\operatorname{Im}\langle E_xE_y^*\rangle=S_0+S_3.
\end{align*}
因此测得$(S_0+S_3)/2$。图\ref{fig:o4analysisa}可用圆基$(1,i)/\sqrt2$、$(1,-i)/\sqrt2$逐个代入：前者经过波片成为水平线偏振，后者成为竖直线偏振，水平片分别全透与阻断。

快轴改为135度，或交换快慢轴，可测得$(S_0-S_3)/2$。波片本身是无损相位元件，输出总强度仍为$S_0$；圆偏振的选择发生在随后的检偏器中。描述“正圆”时必须同时保留时间因子与$S_3$符号约定，避免只用旋向名称造成歧义。
\end{solution}

\begin{example}\label{ex:o4ellipseinverse}
完全偏振光的长短半轴为$A,B$，$A\geq B\geq0$，长轴与$x$成$\psi$，从$+x$向$+y$的旋向取正。求斯托克斯量与椭圆几何的关系，说明哪一个量携带旋向。
\Diagram{fig:o4inverseellipseq}{椭圆的尺寸、倾角和旋向是独立几何信息}{\draw[axis](-2.2,0)--(2.5,0)node[right]{$x$};\draw[axis](0,-1.5)--(0,1.8)node[above]{$y$};\draw[curve,rotate=25](0,0)ellipse(1.8 and .8);\draw[guide](-1.8,-.84)--(1.8,.84);\node at(1,.2){$\psi$};}
\end{example}
\begin{solution}
在图\ref{fig:o4inverseellipseq}的主轴基中，正旋Jones矢量可取$(A,iB)$，所以$S_0=A^2+B^2$、$S_1=A^2-B^2$、$S_2=0$、$S_3=2AB$。将坐标旋转$\psi$，线偏振分量在斯托克斯平面旋转双角，得到
\[
 S_1=(A^2-B^2)\cos2\psi,\quad
 S_2=(A^2-B^2)\sin2\psi,\quad S_3=\pm2AB.
\]
\Diagram{fig:o4inverseellipsea}{相同轨迹反向走动，只改变$S_3$符号}{\begin{scope}[xshift=-2cm]\draw[curve](0,0)ellipse(1.2 and .65);\draw[->](1.2,0)arc(0:65:1.2 and .65);\node[below]at(0,-1){$S_3>0$};\end{scope}\begin{scope}[xshift=2cm]\draw[curve](0,0)ellipse(1.2 and .65);\draw[->](1.2,0)arc(0:-65:1.2 and .65);\node[below]at(0,-1){$S_3<0$};\end{scope}}
图\ref{fig:o4inverseellipsea}相同长短轴和倾角下，旋向只由$S_3$的符号区分。验证$S_1^2+S_2^2+S_3^2=(A^2+B^2)^2$，正是完全偏振条件；圆偏振$A=B$时$S_1=S_2=0$，长轴方向不再可定义，此时$\tan2\psi$不再确定唯一倾角。
\end{solution}

\section{晶体中的波矢与能流}
\begin{example}\label{ex:o4walkoff}
单轴晶体$n_o=1.60,n_e=1.50$，非普通光波矢与光轴成45度。求有效相折射率、能流方向和走离角。忽略频率色散。
\Diagram{fig:o4walkoffq}{波矢与光轴成$45^\circ$，能流偏向远离光轴的一侧}{
\draw[axis](0,0)--(3.8,0)node[right]{$x$};\draw[axis](0,0)--(0,3.4)node[above]{光轴$z$};
\draw[axis](0,0)--(2.15,2.15)node[above left]{$\vct k$};
\draw[axis](0,0)--({2.15*2.56/2.25},2.15)node[right]{$\vct S$};
\draw(0,.85)arc(90:45:.85);\node at(.6,1.12){$\vartheta$};
\draw(45:2.55)arc(45:41.308:2.55);
\draw[thin](43.154:2.55)--(3.1,1.1)node[right]{$\rho$};}
\end{example}
\begin{solution}
图\ref{fig:o4walkoffq}取光轴为$z$、波矢位于$x,z$平面。非普通波色散为$k_x^2/n_e^2+k_z^2/n_o^2=(\omega/c)^2$。令$k=n_{\rm eff}\omega/c$、$k_x=k\sin\vartheta$、$k_z=k\cos\vartheta$，得
\[
 \frac1{n_{\rm eff}^2}=\frac{\sin^2\vartheta}{n_e^2}
 +\frac{\cos^2\vartheta}{n_o^2}.
\]
取$\vartheta=45^\circ$、$n_o=1.60,n_e=1.50$，得到$n_{\rm eff}=1.5476$。这决定沿$\vct k$的相位速度$c/n_{\rm eff}$。
\Diagram{fig:o4walkoffa}{等频椭圆的外法线给出能流方向，虚线延长波矢方向}{
\draw[axis](0,0)--(3.8,0)node[right]{$k_x$};\draw[axis](0,0)--(0,3.5)node[above]{$k_z$（光轴）};
\draw[curve](2.4,0)arc(0:90:2.4 and 2.56);
\pgfmathsetmacro{\p}{1/sqrt(1/2.4^2+1/2.56^2)}
\fill(\p,\p)circle(1.5pt)node[below right]{$P$};
\draw[axis](0,0)--(\p,\p)node[midway,below right]{$\vct k$};
\draw[guide](\p,\p)--({\p+1.05},{\p+1.05});
\draw[axis](\p,\p)--({\p+.95*2.56/2.25},{\p+.95})node[right]{$\vct S$};
\draw(0,.7)arc(90:45:.7);\node at(.4,.94){$45^\circ$};
\node[align=left]at(3.4,.65){$\vartheta_S=48.69^\circ$\\$\rho=3.69^\circ$};}
无色散时，群速度为$\vct v_g=\nabla_{\vct k}\omega$。固定$k_z$，对色散式的两边关于$k_x$求导；再固定$k_x$，关于$k_z$求导，分别有
\[
 \frac{2k_x}{n_e^2}=\frac{2\omega}{c^2}\frac{\partial\omega}{\partial k_x},
 \qquad
 \frac{2k_z}{n_o^2}=\frac{2\omega}{c^2}\frac{\partial\omega}{\partial k_z}.
\]
解出两个偏导数，得到
\[
 v_{gx}=\frac{c^2k_x}{\omega n_e^2},\qquad
 v_{gz}=\frac{c^2k_z}{\omega n_o^2},
\]
因此图\ref{fig:o4walkoffa}的能流方向满足
\[
 \tan\vartheta_S=\frac{v_{gx}}{v_{gz}}
 =\frac{n_o^2}{n_e^2}\tan\vartheta
 =\frac{2.56}{2.25},\quad \vartheta_S=48.69^\circ.
\]
走离角$\rho=\vartheta_S-\vartheta=3.69^\circ$，方向远离光轴。图中的箭头沿等频椭圆在$P$点的外法线，虚线则沿波矢方向延长，两者夹角为$\rho$。若入射点取原点、出射端面为$z=h$，光束中心横坐标为$h\tan\vartheta_S$；相位法线方向对应$h\tan\vartheta$，二者横向差为$h(\tan\vartheta_S-\tan\vartheta)$。因此光斑位置与波面相位不能用同一个折射角描述。
\end{solution}

\begin{example}\label{ex:o4reflectionpolar}
自然光射向无损介质界面。给定菲涅耳反射率$R_s,R_p$，求反射光的线偏振度和总强度，讨论正入射、布儒斯特角和掠入射三个极限。
\Diagram{fig:o4reflectionq}{自然光的两正交份额分别被不同反射率筛选}{\node at(0,.6){$I_0/2:s$};\node at(0,-.6){$I_0/2:p$};\draw[->](1.2,.6)--(2.6,.6)node[midway,above]{$R_s$};\draw[->](1.2,-.6)--(2.6,-.6)node[midway,below]{$R_p$};\node[right]at(2.8,0){反射束};}
\end{example}
\begin{solution}
图\ref{fig:o4reflectionq}反射强度$I_r=I_0(R_s+R_p)/2$，两正交分量不相干，所以线偏振度$\mathcal P=|R_s-R_p|/(R_s+R_p)$。若$R_s>R_p$，主要偏振沿$s$。
\Diagram{fig:o4reflectiona}{布儒斯特角处纯$s$，两端极限偏振度回落}{\draw[axis](0,0)--(5.2,0)node[right]{$i$};\draw[axis](0,0)--(0,2.8)node[above]{$\mathcal P$};\draw[curve](0,0)..controls(1.7,.2)and(2.6,2.3)..(3.1,2.3)..controls(3.8,2.3)and(4.4,.7)..(4.8,0);\node[below]at(3.1,0){$i_B$};\node[below]at(4.8,0){$\pi/2$};}
正入射$R_s=R_p$，偏振度0；布儒斯特角$R_p=0$，偏振度1；掠入射两者都趋1，偏振度又趋0，图\ref{fig:o4reflectiona}表示这种趋势。若总反射强度为零，例如完全无折射率差的界面，则比值为$0/0$，无反射光可谈偏振，不能随意给偏振度0或1。
\end{solution}

\section{全反射延迟与边界附近能流}
\begin{example}\label{ex:o4rhomb}
玻璃向空气全反射两次，两次入射面共面且偏振基连续对应。求实现总相对延迟$\pi/2$所需的单次相位差条件；对$n=1.5$推导可用入射角方程。
\Diagram{fig:o4rhombq}{两次全反射的相位差相加}{\draw[curve](0,0)--(4,0)--(5,2)--(1,2)--cycle;\draw[->](.4,.8)--(1.2,0)--(3.2,2)--(4.4,.8);\node[below]at(1.2,0){$R_1$};\node[above]at(3.2,2){$R_2$};}
\end{example}
\begin{solution}
设$m=n_2/n_1=2/3$，$g=\sqrt{\sin^2i-m^2}$。单次相位$\phi_s=-2\arctan(g/\cos i)$、$\phi_p=-2\arctan(g/(m^2\cos i))$。图\ref{fig:o4rhombq}两次基底对应相同，需$|\phi_p-\phi_s|=\pi/4$。
\Diagram{fig:o4rhomba}{两相量相位差每次增加45度，两次累积90度}{\draw[axis](0,0)--(2.8,0)node[right]{$s$参考};\draw[->](0,0)--(1.6,-1.6)node[right]{一次};\draw[->,thick](0,0)--(0,-2.3)node[below]{两次};\draw[guide](2.3,0)arc(0:-90:2.3);}
令$x=g/\cos i>0$，用反正切差公式，条件为
\[
 \frac{x(1/m^2-1)}{1+x^2/m^2}=\tan(\pi/8).
\]
记$u=\tan(\pi/8)$，化为二次式$ux^2-(1-m^2)x+m^2u=0$，两根为
\[
 x_\pm=\frac{1-m^2\pm\sqrt{(1-m^2)^2-4m^2u^2}}{2u}.
\]
代入$m=2/3$，得$x_-=0.59795$、$x_+=0.74328$。又由$x^2=(\sin^2i-m^2)/(1-\sin^2i)$，整理为$(1+x^2)\sin^2i=x^2+m^2$，故
\[
 i_\pm=\arcsin\sqrt{\frac{x_\pm^2+m^2}{1+x_\pm^2}}.
\]
两解约为50.23度和53.26度，均大于临界角41.81度。图\ref{fig:o4rhomba}画的是相对延迟的累积；实际棱镜的入射、出射面若还引入振幅失配或相位差，也应计入总Jones矩阵。
\end{solution}

\begin{example}\label{ex:o4evanescentflux}
$s$偏振倏逝波在$z>0$介质中为$E_y=E_t\ee^{ik_xx-\kappa z-i\omega t}$，$k_x,\kappa$均为正实数。直接由麦克斯韦方程求平均坡印廷矢量，验证法向能流为零而沿界面不为零。
\Diagram{fig:o4evanq}{倏逝振幅沿法向衰减，相位沿界面前进}{\draw[axis](-.3,0)--(4.5,0)node[right]{$x$};\draw[axis](0,0)--(0,3.2)node[above]{$z$};\foreach\y/\l in{.3/3.8,.8/2.6,1.3/1.8,1.8/1.2}{\draw[curve](.2,\y)--(\l,\y);}\node[right]at(2,2.5){$\ee^{-\kappa z}$};}
\end{example}
\begin{solution}
取复波矢$\vct k=(k_x,0,i\kappa)$，图\ref{fig:o4evanq}的$\vct H=(\vct k\times\vct E)/(\mu\omega)$，所以$H_x=-i\kappa E_y/(\mu\omega)$、$H_z=k_xE_y/(\mu\omega)$。复坡印廷乘积为
\[
 \vct E\times\vct H^*
 =\frac{|E_y|^2}{\mu\omega}(k_x\hat{\vct x}-i\kappa\hat{\vct z}).
\]
\Diagram{fig:o4evana}{时间平均实能流平行界面，法向虚部表示往返储能交换}{\draw[curve](-.5,0)--(4.5,0);\foreach\y in{.5,1,1.5}{\draw[axis](.3,\y)--({3.5*exp(-.5*\y)},\y);}\node[right]at(3,1.7){$\langle S_z\rangle=0$};}
取一半实部，得$\langle S_x\rangle=k_x|E_t|^2\ee^{-2\kappa z}/(2\mu\omega)$，$\langle S_z\rangle=0$，如图\ref{fig:o4evana}。场在低折射率侧非零与无净法向透射并不矛盾。若存在第二界面产生另一指数支，交叉项可给实法向能流，必须重新计算而不能保留单支结果。
\end{solution}

\section{偏振元件的损耗与色散}
\subsection{相位延迟和选择性衰减改变不同的量}
以下继续以$C$表示相干矩阵、$J$表示Jones矩阵，元件作用为$C'=JCJ^\dagger$。第\ref{th:o4coherence}节已证明偏振度由$C$的迹与行列式确定；无损波片保持这两个量。二向色元件则对两个本征方向作不同衰减，改变$C$的本征值比例，因此可能以损失通光量为代价提高偏振度。

设主轴方向上的能量透过率为$t_x,t_y$，无额外相位延迟，则$D=\operatorname{diag}(\sqrt{t_x},\sqrt{t_y})$。若自然光入射，$C=I_0\mathbb I/2$，则$C_D=DCD^\dagger=I_0\operatorname{diag}(t_x,t_y)/2$。取其迹与两个本征值之差，得到出射强度和偏振度
\[
 I=\frac{I_0(t_x+t_y)}2,\qquad
 P=\frac{|t_x-t_y|}{t_x+t_y}.
\]
以上偏振度公式要求$I_0>0$且$t_x+t_y>0$；若两方向都完全遮光，出射强度为零，偏振度无定义。随后接一块无损波片，强度和偏振度不变，偏振成分却可变为椭圆偏振。两元件的先后顺序须按$QD$或$DQ$计算，它们一般不对易。

\begin{example}部分圆偏振光的产生与检偏
自然光强度为$I_0$，先通过主轴能量透过率$t_x=1,t_y=1/4$的二向色元件，再通过快轴与$x$成$45^\circ$的理想四分之一波片。按本讲Jones矩阵约定，删去共同相位后，$x$线偏振变为$(1,-\mathrm i)^{\mathsf T}/\sqrt2$。求出射强度、偏振度及旋转线检偏器的强度范围。比较交换两元件顺序的结果。
\ExpansionDiagram{fig:dep-o4-loss-q}{先选择性衰减、后相位延迟的偏振光路}{
\draw[->](0,0)--(7,0);\draw[fill=white](1,-.7)rectangle(2,.7)node[midway]{$D$};\draw[fill=white](3,-.7)rectangle(4,.7)node[midway]{$Q$};\draw[fill=white](5,-.7)rectangle(6,.7)node[midway]{$A_\theta$};\node[above]at(1.5,.9){$t_x\ne t_y$};\node[above]at(3.5,.9){$45^\circ$};}
\end{example}
\begin{solution}
由图\ref{fig:dep-o4-loss-q}先经$D$，得$C_D=I_0\operatorname{diag}(1,1/4)/2$。按第\ref{th:o4coherence}节的本征值分解，将它分成自然光与完全偏振光：
\[
 C_D=\frac{I_0}{8}\mathbb I+\frac{3I_0}{8}
 \begin{pmatrix}1&0\\0&0\end{pmatrix}.
\]
第一项是总强度$I_0/4$的自然光，第二项是强度$3I_0/8$的$x$线偏振。通过$Q$后，第一项不变，第二项按题设变为圆偏振，所以
\[
 I_{\rm out}=\frac{5I_0}{8},\qquad
 P=\frac{3I_0/8}{5I_0/8}=\frac35.
\]
\ExpansionDiagram{fig:dep-o4-loss-a}{相干矩阵分解中的自然成分与完全偏振成分}{
\draw(0,0)rectangle(2,1.2);\node at(1,.8){自然光};\node at(1,.35){$I_0/4$};\node at(2.5,.6){$+$};\draw(3,0)rectangle(5,1.2);\node at(4,.8){圆偏振};\node at(4,.35){$3I_0/8$};\node at(5.5,.6){$\longrightarrow$};\node at(7,.6){$P=3/5$};}
图\ref{fig:dep-o4-loss-a}中两种成分通过任意方向的理想线检偏器，都各透过一半。因此$I(\theta)=5I_0/16$恒定，虽然$P=3/5$并不为零。线检偏扫描没有起伏不能单独证明自然光。

若交换为先$Q$后$D$，自然光先经无损波片仍是自然光，最后输出恰为$C_D$。其强度和$P$仍相同，但偏振成分是线偏振，检偏强度为
\[
 I(\theta)=\frac{I_0}{2}\cos^2\theta+\frac{I_0}{8}\sin^2\theta,
\]
最大$I_0/2$、最小$I_0/8$。顺序变化改变了相干矩阵的方向结构，而本题两个顺序恰好有相同特征值；这不意味着任意入射态都保持相同输出强度。
\end{solution}

\subsection{复合波片的一阶消色差条件}
单片延迟量$\delta=2\pi h\Delta n(\lambda)/\lambda$同时受到折射率色散和$1/\lambda$因子的影响。即使忽略材料色散，延迟量也不与波长无关。两材料波片主轴平行或互相垂直时，延迟量可以代数相加：
\[
 \delta(\lambda)=\frac{2\pi}{\lambda}
 [h_1\Delta n_1(\lambda)+s h_2\Delta n_2(\lambda)],\qquad s=\pm1.
\]
这里$s=-1$表示把第二片的快慢轴对调。若两片轴任意夹角，应乘Jones矩阵，不能再直接加延迟量。

令$L(\lambda)=h_1\Delta n_1+s h_2\Delta n_2$。厚度固定时，对$\delta=2\pi L/\lambda$用商法则求导，得
\[
 \delta'(\lambda)=2\pi\left[\frac{L'(\lambda)}{\lambda}-\frac{L(\lambda)}{\lambda^2}\right].
\]
在目标波长$\lambda_0$处设计零级四分之一波片，并令$\delta'(\lambda_0)=0$，需要
\[
 L(\lambda_0)=\frac{\lambda_0}{4},\qquad
 L'(\lambda_0)=\frac{L(\lambda_0)}{\lambda_0}=\frac14.
\]
第二式中的非零$L'$补偿了$1/\lambda$因子的变化。厚度误差、温度膨胀和材料高阶色散仍须计入延迟量；一阶消色差条件没有消除这些影响。

\begin{example}两材料波片的厚度设计
两材料在$\lambda_0=500\,\mathrm{nm}$附近满足$\Delta n_1=0.010-10^{-5}(\lambda-500)$、$\Delta n_2=0.020-5\times10^{-5}(\lambda-500)$，式中波长数值以纳米为单位。将两片主轴交叉，求使总延迟为$\pi/2$且一阶消色差的正厚度。
\ExpansionDiagram{fig:dep-o4-achrom-q}{第二片主轴对调，使两片延迟量相减}{
\draw(0,0)rectangle(2,1.8);\draw[<->](.3,.9)--(1.7,.9)node[midway,above]{快轴1};\draw(3,0)rectangle(5,1.8);\draw[<->](4,.2)--(4,1.6);\node[right]at(4,1){快轴2};\node[below]at(1,0){$h_1$};\node[below]at(4,0){$h_2$};}
\end{example}
\begin{solution}
按图\ref{fig:dep-o4-achrom-q}取$s=-1$。厚度也用纳米表示，中心延迟条件和导数条件分别为
\[
 0.010h_1-0.020h_2=125,
 \qquad -10^{-5}h_1+5\times10^{-5}h_2=\frac14.
\]
第一式给$h_1=12500+2h_2$，代入第二式为$-0.125+3\times10^{-5}h_2=0.25$，故
\[
 h_2=12500\,\mathrm{nm}=12.5\,\mu\mathrm m,
 \qquad h_1=37.5\,\mu\mathrm m.
\]
两片单独的光程差为375 nm和250 nm，其差125 nm给$\pi/2$；导数之差为$-0.375-(-0.625)=0.25$，同时满足要求。
\ExpansionDiagram{fig:dep-o4-achrom-a}{两片光程差的差值随波长保持所需的一阶斜率}{
\draw[->](0,0)--(6,0)node[right]{$\lambda$};\draw[->](0,0)--(0,3.1)node[above]{光程差};\draw[thick](.5,2.3)--(5.3,1.7)node[right]{$h_1\Delta n_1$};\draw[dashed](.5,1.6)--(5.3,.6)node[right]{$h_2\Delta n_2$};\draw[<->](3,1.08)--(3,1.99)node[midway,right]{$L(\lambda)$};}
图\ref{fig:dep-o4-achrom-a}中的两条线都下降，但第二条下降更快，所以差值$L$随波长增加。本题给定的严格线性模型恰使$L=\lambda/4$，在该模型适用区间内总延迟恒定；实际材料的二阶以上色散会留下残差。这里不能凭线性模型推断无限宽带性能。
\end{solution}

\section{闭合偏振路径中的干涉相位}
\subsection{单束可以删去的共同相位，在两臂间仍可测量}
两个非零Jones矢量若只相差共同相位，其单束强度和偏振椭圆相同，与同一参考光干涉时却会产生不同的条纹位置。依次投影到几个偏振态后，即使偏振方向回到初态，也可能留下可测的路径相位。因此比较两臂时，应保留各元件作用后的完整复振幅。

归一化偏振态$|a\rangle$的理想投影元件写$P_a=|a\rangle\langle a|$。它选择该态并保留投影的复数系数。初态$|a\rangle$依次经$P_b,P_c,P_a$后，输出为
\[
 |a\rangle\langle a|c\rangle\langle c|b\rangle\langle b|a\rangle.
\]
三个内积的积包含透过幅度与路径相位。把任一态改写为$\ee^{\mathrm i\alpha}|a\rangle$并不改变投影算符，积中的相位也成对抵消。因此所求相位不是任选Jones代表元造成的假象。实际元件还可能有传播相位，干涉装置中需事先平衡或把它们明确加入模型。

\begin{example}交换偏振投影顺序，最后方向相同而条纹移动
取$|H\rangle=(1,0)^{\mathsf T}$、$|D\rangle=(1,1)^{\mathsf T}/\sqrt2$、$|R\rangle=(1,\mathrm i)^{\mathsf T}/\sqrt2$。干涉仪两臂在进入元件前均为场$E_0|H\rangle$。参考臂附加相位为$\delta$，信号臂先经$P_D$、再$P_R$、最后$P_H$。求信号透过率与干涉强度；再交换$P_D,P_R$，比较条纹。投影元件按上述算符定义，其他光程相位已补偿。
\ExpansionDiagram{fig:dep-o4-loop-q}{两条信号路径最终都回到水平偏振，参考臂保留相位基准}{
\node[draw]at(0,1){$H$};\node[draw]at(2,1){$D$};\node[draw]at(4,1){$R$};\node[draw]at(6,1){$H$};\draw[->](.4,1)--(1.6,1);\draw[->](2.4,1)--(3.6,1);\draw[->](4.4,1)--(5.6,1);\draw[->](0,-.3)--(6,-.3)node[midway,below]{参考臂$\ee^{\mathrm i\delta}H$};}
\end{example}
\begin{solution}
图\ref{fig:dep-o4-loop-q}第一步幅度系数为$\langle D|H\rangle=1/\sqrt2$，第二步为$\langle R|D\rangle=(1-\mathrm i)/2$，第三步为$\langle H|R\rangle=1/\sqrt2$。保留三个复系数相乘，得
\[
 c=\frac{1-\mathrm i}{4}=\frac1{2\sqrt2}\ee^{-\mathrm i\pi/4}.
\]
信号透过率$|c|^2=1/8$，最后偏振仍为$H$。参考臂与信号臂叠加时
\[
 I=I_0|\ee^{\mathrm i\delta}+c|^2
 =I_0\left[\frac98+\frac1{\sqrt2}\cos(\delta+\pi/4)\right],
\]
其中$I_0=|E_0|^2$采用同一强度单位。
\ExpansionDiagram{fig:dep-o4-loop-a}{两种顺序的信号系数互为共轭，参考相位扫描看到相反偏移}{
\draw[->](-.2,0)--(3.5,0)node[right]{实部};\draw[->](0,-2)--(0,2)node[above]{虚部};\draw[->,thick](0,0)--(2,-2)node[right]{$(1-\mathrm i)/4$};\draw[->,thick](0,0)--(2,2)node[right]{$(1+\mathrm i)/4$};\draw[dashed](2,-2)--(2,2);}
交换后$c'=\langle H|D\rangle\langle D|R\rangle\langle R|H\rangle=(1+\mathrm i)/4$。图\ref{fig:dep-o4-loop-a}中两信号模相同、辐角相差$\pi/2$，新强度为$I_0[9/8+\cos(\delta-\pi/4)/\sqrt2]$。两条纹相差四分之一个周期，可见度同为$4\sqrt2/9$。

两种顺序得到的信号具有相同强度和偏振，与共同参考臂叠加后才显出路径相位差。仅计算单束偏振时可以删去共同相位；计算干涉时，须另行保留它相对于参考臂的值。逐片归一化并丢掉投影复系数，会同时丢失透过率和条纹位移。
\end{solution}

\section{习题与解答}
\begin{exercise}\label{ex:o4twohalf}
两块理想半波片的快轴方向依次为$\alpha_1,\alpha_2$。求合成Jones矩阵和线偏振转角，说明为何要保留本征轴方向而不能只给光场乘一个共同相位。
\Diagram{fig:o4twohalfq}{两半波片的先后轴方向决定合成旋转}{\draw[axis](-1,0)--(5.5,0);\draw[curve](0,-1)--(0,1);\draw[curve](3,-1)--(3,1);\node[above]at(0,1){$H(\alpha_1)$};\node[above]at(3,1){$H(\alpha_2)$};}
\end{exercise}
\begin{solution}
删去每片的共同相位，图\ref{fig:o4twohalfq}半波片矩阵为
\[
 H(\alpha)=\begin{pmatrix}\cos2\alpha&\sin2\alpha\\
 \sin2\alpha&-\cos2\alpha\end{pmatrix}.
\]
令$c_j=\cos2\alpha_j$、$s_j=\sin2\alpha_j$，按传播顺序相乘得到
\[
 H(\alpha_2)H(\alpha_1)=
 \begin{pmatrix}c_2c_1+s_2s_1&c_2s_1-s_2c_1\\
 s_2c_1-c_2s_1&s_2s_1+c_2c_1\end{pmatrix}
 =R[2(\alpha_2-\alpha_1)].
\]
其中$R(\theta)=\bigl(\begin{smallmatrix}\cos\theta&-\sin\theta\\\sin\theta&\cos\theta\end{smallmatrix}\bigr)$。因此任意线偏振方向转过$2(\alpha_2-\alpha_1)$，椭圆的长轴同样旋转，旋向经两次翻转后恢复。
\Diagram{fig:o4twohalfa}{图取$\alpha_1=20^\circ,\alpha_2=60^\circ$，合成旋转$80^\circ$}{\draw[axis](0,0)--(3,0)node[right]{原方向};\draw[guide](0,0)--(2.53717,.92345)node[right]{$\alpha_1$};\draw[guide](0,0)--(1.5,2.59808)node[above]{$\alpha_2$};\draw[axis](0,0)--(.45149,2.56050)node[above left]{输出};\draw[->](1,0)arc(0:80:1);\node[below]at(1,-.6){$2(\alpha_2-\alpha_1)$};}
图\ref{fig:o4twohalfa}反映两次镜像的几何关系。交换波片顺序会使旋转角反号；快轴平行时合成为单位矩阵。半波片的相对延迟为$\pi$，不是对所有输入都加同一个标量相位；忽略本征轴会丢失整个旋转效应。
\end{solution}

\begin{exercise}\label{ex:o4geometricphase}
干涉仪两臂输入同一正圆偏振态$(1,i)^{\mathsf T}/\sqrt2$。一臂无偏振元件，另一臂依次通过快轴0度和$\alpha$的两块半波片。普通传播相位已补偿，两臂等强。求输出偏振、相对相位和合束强度。
\Diagram{fig:o4geometricq}{两臂最终偏振相同，波片组合仍可改变干涉相位}{\draw[->](0,1)--(5.5,1);\draw[->](0,-.5)--(5.5,-.5);\draw[curve](1.5,-1)--(1.5,0);\draw[curve](3.5,-1)--(3.5,0);\node[below]at(1.5,-1){$H(0)$};\node[below]at(3.5,-1){$H(\alpha)$};\node[right]at(5.5,.3){合束};}
\end{exercise}
\begin{solution}
由上一题，图\ref{fig:o4geometricq}测试臂的矩阵为$R(2\alpha)$。对圆Jones矢量直接作用，
\[
 R(2\alpha)\frac1{\sqrt2}\begin{pmatrix}1\\i\end{pmatrix}
 =\ee^{-2i\alpha}\frac1{\sqrt2}\begin{pmatrix}1\\i\end{pmatrix}.
\]
所以输出仍是同一正圆偏振，公共相位相对参考臂却改变$-2\alpha$。若单臂强度为$I_1$，两臂合场比例$1+\ee^{-2i\alpha}$，强度为$4I_1\cos^2\alpha$。
\Diagram{fig:o4geometrica}{两臂振幅相等，圆态的公共相位可由干涉测出}{\draw[axis](0,0)--(2.5,0)node[right]{参考相量};\draw[axis](0,0)--(1,-2.291288)node[right]{测试相量};\draw[guide](2.5,0)arc(0:-66.42182:2.5);\node at(1.4,-.5){$-2\alpha$};}
图\ref{fig:o4geometrica}的相位随波片转角改变，$\alpha=0$相长，$\alpha=\pi/2$相消。单臂的偏振椭圆保持不变，两臂的干涉强度却随这个公共相位变化。计算中删去的两片材料本身公共传播相位必须确已补偿，本题只隔离随轴角变化的部分。
\end{solution}

\begin{exercise}\label{ex:o4compensator}
完全偏振光Jones矢量为$(1,2\ee^{i\pi/3})^{\mathsf T}$。只允许转动一块四分之一波片，求可使出射成为线偏振的轴方向，并给出线偏振输出与椭圆主轴的关系。
\Diagram{fig:o4compq}{先由椭圆主轴确定补偿波片的本征轴}{\draw[axis](-1.8,0)--(2,0)node[right]{$x$};\draw[axis](0,-2.2)--(0,2.4)node[above]{$y$};\draw[curve,domain=0:360,samples=100]plot({cos(\x)},{2*cos(\x-60)});\node[right]at(1.5,1.7){待补偿};}
\end{exercise}
\begin{solution}
图\ref{fig:o4compq}有$a=1,b=2,\delta=\pi/3$。主轴角满足$\tan2\psi=-2/3$，用$\operatorname{atan2}(2,-3)$选长轴，$\psi\simeq73.155^\circ$。长短半轴平方为$A^2=(5+\sqrt{13})/2$、$B^2=(5-\sqrt{13})/2$。
\Diagram{fig:o4compa}{快慢轴对准长短轴后，90度相差被补成0或180度}{\draw[axis](-2,0)--(2.4,0)node[right]{长轴$\xi$};\draw[axis](0,-1.6)--(0,2)node[above]{短轴$\eta$};\draw[curve](0,0)ellipse(1.8 and .725);\draw[very thick](-1.8,.725)--(1.8,-.725);\draw[guide](-1.8,-.725)--(1.8,.725);}
原态$S_3>0$，在主轴基删公共相位可写$(A,iB)$。若快轴沿长轴、慢轴沿短轴，波片后为$(A,-B)$，输出线方向为$\psi-\arctan(B/A)$；若快慢互换，成为$(iA,iB)$，方向为$\psi+\arctan(B/A)$。图\ref{fig:o4compa}画出两种补偿线。因此波片快轴可沿$73.155^\circ$或其正交方向，输出方向由补偿后的两个非零分量合成，一般偏离原椭圆长轴。
\end{solution}

\begin{exercise}\label{ex:o4ellipsereconstruct}
完全偏振光经旋转线检偏器的最大、最小强度比为9:1，最大方向与$x$成30度。求椭圆长短轴比及长轴方向；仅凭这些数据能否确定旋向？
\Diagram{fig:o4reconstructq}{检偏最大方向与椭圆长轴一致}{\draw[axis](-2,0)--(2.5,0)node[right]{$x$};\draw[curve,rotate=30](0,0)ellipse(1.8 and .6);\draw[guide](-1.8,-1.039)--(1.8,1.039);\node[right]at(1.8,1.039){30度};}
\end{exercise}
\begin{solution}
图\ref{fig:o4reconstructq}完全偏振条件下，沿主轴检偏得到强度分别正比$A^2,B^2$，所以$A/B=3$，长轴为30度模180度。强度比不是半轴比9。
\Diagram{fig:o4reconstructa}{同一椭圆轨迹允许两种相反旋向}{\draw[curve](0,0)ellipse(1.6 and .6);\draw[->](1.6,0)arc(0:60:1.6 and .6);\draw[->](1.6,0)arc(0:-60:1.6 and .6);\node[right]at(2,.3){需圆偏振分析};}
图\ref{fig:o4reconstructa}两种旋向有相同线检偏曲线，需波片测$S_3$符号。若没有“完全偏振”前提，9:1还可能来自自然光与线偏振混合，不能唯一恢复椭圆半轴。
\end{solution}

\begin{exercise}\label{ex:o4halfcircular}
证明理想半波片将任意完全偏振光的椭圆旋向反转，并求其长轴方向变化。输入长轴角$\psi$，快轴角$\alpha$。
\Diagram{fig:o4halfcircq}{半波片在快慢轴上把一个分量反号}{\draw[axis](-2,0)--(2.3,0)node[right]{快轴};\draw[axis](0,-1.7)--(0,2)node[above]{慢轴};\draw[curve,rotate=25](0,0)ellipse(1.6 and .7);}
\end{exercise}
\begin{solution}
图\ref{fig:o4halfcircq}中，在波片快、慢轴基底去掉共同相位，半波片为$H=\operatorname{diag}(1,-1)$。对任意时间的实电场，它使$(E_f(t),E_s(t))$变为$(E_f(t),-E_s(t))$，即关于快轴的镜像。镜像保留长短半轴，长轴相对快轴的夹角由$\psi-\alpha$变为$-(\psi-\alpha)$，所以非圆椭圆的输出方向为
\[
 \psi_o=2\alpha-\psi\pmod\pi.
\]
\Diagram{fig:o4halfcirca}{镜像变换保留椭圆形状，反转沿轨迹的有向旋转}{\draw[axis](-2.4,0)--(2.6,0)node[right]{快轴};\draw[curve,rotate=25](0,0)ellipse(1.8 and .65);\draw[guide,rotate=-25](0,0)ellipse(1.8 and .65);\node[below]at(0,-1.8){$\psi_o=2\alpha-\psi$};}
图\ref{fig:o4halfcirca}中的轨迹形状不变，有向面积却反号。代数上，轴基底的互相关$E_fE_s^*$变为$-E_fE_s^*$，故$S_3$变号；绕传播轴旋转实坐标基底不改变$S_3$，这一结论在实验室基底同样成立。完全偏振光的椭圆率大小保持，旋向反转。

有两个退化情形须分开：线偏振的$S_3=0$，本来没有椭圆旋向，只需使用线方向公式；圆偏振满足长短半轴相同，长轴方向$\psi$没有定义，不能再用长轴角描述输出。对圆Jones矢量直接作用$H$，$(1,i)$变为$(1,-i)$，能明确判断圆旋向反转。公共相位改变时间参数的起点，不改变上述几何结论。
\end{solution}

\begin{exercise}\label{ex:o4twoquarter}
两块相同四分之一波片快轴平行或正交。分别求合成Jones矩阵，并说明两片都存在时偏振是否一定改变。
\Diagram{fig:o4twoquarterq}{两片波片的相对轴角决定相位差相加或相消}{\draw[curve](0,-1)--(0,1);\draw[curve](3,-1)--(3,1);\draw[axis](-1,0)--(4.5,0);\node[above]at(0,1){$Q_1$};\node[above]at(3,1){$Q_2$};}
\end{exercise}
\begin{solution}
按图\ref{fig:o4twoquarterq}，快轴平行时两矩阵均为$\operatorname{diag}(1,i)$，乘积$\operatorname{diag}(1,-1)$，等效半波片。快轴正交时第二片在同一基底为$\operatorname{diag}(i,1)$，乘积$iI$。
\Diagram{fig:o4twoquartera}{正交轴相消后只剩共同相位，所有输入偏振态保持}{\node at(0,.5){平行：$Q^2=H$};\node at(0,-.6){正交：$Q_\perp Q=iI$};\draw[->](2,.5)--(3.3,.5);\draw[->](2,-.6)--(3.3,-.6);\node[right]at(3.5,.5){相差累加};\node[right]at(3.5,-.6){偏振不变};}
图\ref{fig:o4twoquartera}的正交轴组合在所用约化矩阵中为$iI$，所有输入偏振态保持。若恢复每片删去的快轴传播相位$\phi_{f1},\phi_{f2}$，完整矩阵为$\ee^{i(\phi_{f1}+\phi_{f2})}iI$；与参考臂干涉时须保留这整个共同相位，不能仅由$iI$断言实物插入后的总相移恰为$\pi/2$。平行轴组合虽等效半波片，对沿其本征轴的输入仍只改变共同相位。
\end{solution}

\begin{exercise}\label{ex:o4criticaldecay}
全反射入射角为$i=i_c+\epsilon$，$0<\epsilon\ll1$，$i_c=\arcsin(n_2/n_1)$。求倏逝振幅衰减长度随$\epsilon$的领先阶，说明临界角附近的尺度。
\Diagram{fig:o4criticalq}{刚超过临界角，低折射率侧场衰减很慢}{\draw[curve](-2.3,0)--(2.3,0);\draw[guide](0,-1.5)--(0,2);\draw[->](-1.6,-1.3)--(0,0);\draw[->](0,0)--(1.6,-1.3);\node[right]at(1.2,1){长衰减尾部};}
\end{exercise}
\begin{solution}
图\ref{fig:o4criticalq}的$\kappa=k_0\sqrt{n_1^2\sin^2i-n_2^2}$。记$F(i)=n_1^2\sin^2i-n_2^2$，对平方和三角函数逐层求导有$F'(i)=2n_1^2\sin i\cos i$。利用$F(i_c)=0$、$\sin i_c=n_2/n_1$，得到
\begin{align*}
 F(i_c+\epsilon)&=F(i_c)+F'(i_c)\epsilon+O(\epsilon^2)\\
 &=2n_2\sqrt{n_1^2-n_2^2}\,\epsilon+O(\epsilon^2).
\end{align*}
将这一领先项代入$\ell=1/(k_0\sqrt F)$，得
\[
 \ell=\frac1\kappa\simeq
 \frac{\lambda_0}{2\pi\sqrt{2n_2\sqrt{n_1^2-n_2^2}\,\epsilon}}.
\]
\Diagram{fig:o4criticala}{衰减长度在临界角右侧按$\epsilon^{-1/2}$增大}{\draw[axis](0,0)--(4.5,0)node[right]{$\epsilon$};\draw[axis](0,0)--(0,3.5)node[above]{$\ell$};\draw[curve,domain=.2:4,samples=70]plot(\x,{1.4/sqrt(\x)});}
图\ref{fig:o4criticala}的发散说明接近临界角时必须考虑有限光束角谱和第二界面距离。若光束有角度展宽，一部分成分可能低于临界角，单一平面波的无穷衰减长度不再完整描述实际束形。
\end{solution}

\begin{exercise}\label{ex:o4magneticinterface}
正常入射到无损磁性介质界面，介质波阻抗分别为$Z_1=\sqrt{\mu_1/\varepsilon_1}$、$Z_2=\sqrt{\mu_2/\varepsilon_2}$。推导共同切向电场基底的$r,t,R,T$，说明何时可只用折射率。
\Diagram{fig:o4magneticq}{正常入射时电场与磁场的比例由波阻抗决定}{\draw[curve](2,-1.3)--(2,1.3);\draw[->](0,.5)--(1.9,.5);\draw[->](1.9,-.4)--(0,-.4);\draw[->](2.1,.5)--(4.5,.5);\node[above]at(.8,1){$Z_1$};\node[above]at(3.3,1){$Z_2$};}
\end{exercise}
\begin{solution}
图\ref{fig:o4magneticq}中，统一取切向电场沿$x$，入射波沿$+z$、反射波沿$-z$。由$\vct H=\hat{\vct k}\times\vct E/Z$，反射磁场相对其电场多一个方向反号。将入射电场幅度取1，边界条件为
\[
 1+r=t,\qquad \frac{1-r}{Z_1}=\frac t{Z_2}.
\]
用$t=1+r$代入，得到$Z_2-Z_2r=Z_1+Z_1r$，因此$r=(Z_2-Z_1)/(Z_2+Z_1)$、$t=2Z_2/(Z_1+Z_2)$。
\Diagram{fig:o4magnetica}{阻抗匹配可无反射，即使相折射率不同}{\draw[curve](0,-1.5)--(0,1.7);\node[above]at(-1.7,1.5){$Z_1$};\node[above]at(1.8,1.5){$Z_2=Z_1$};\draw[axis](-3,0)--(3.2,0)node[right]{$z$};\draw[->,thick](-2.3,0)--(-2.3,1)node[above]{$E_i$};\draw[->,thick](1.5,0)--(1.5,1)node[above]{$E_t=E_i$};\node[below]at(-1.6,-.6){$E_r=0$};\node[below]at(1.5,-.6){$T=1$};}
无损介质$Z_1,Z_2>0$，平面波平均能流大小为$|E|^2/(2Z)$。图\ref{fig:o4magnetica}中反、入射处于同一介质，透射处于介质2，所以
\[
 R=|r|^2,\qquad T=\frac{Z_1}{Z_2}|t|^2
 =\frac{4Z_1Z_2}{(Z_1+Z_2)^2}.
\]
两者相加的分子为$(Z_2-Z_1)^2+4Z_1Z_2=(Z_1+Z_2)^2$，故$R+T=1$。当$Z_1=Z_2$时无反射，即使两介质的折射率不同也成立，因为阻抗依赖$\mu/\varepsilon$，折射率依赖$\mu\varepsilon$。两介质均非磁性、磁导率相同时，$Z\propto1/n$，才可把阻抗比换成反折射率比，恢复通常正常入射公式。
\end{solution}

\begin{exercise}\label{ex:o4uncertaintystate}
两束互不相干光分别为正、负圆偏振，强度为$3I$和$I$。求总斯托克斯量和偏振度；旋转线检偏器能否观察到强度变化？
\Diagram{fig:o4mixtureq}{两相反圆偏振分量不相干地叠加}{\begin{scope}[xshift=-2cm]\draw[curve](0,0)circle(.8);\draw[->](.8,0)arc(0:70:.8);\node[below]at(0,-1.1){$3I$};\end{scope}\begin{scope}[xshift=2cm]\draw[curve](0,0)circle(.8);\draw[->](.8,0)arc(0:-70:.8);\node[below]at(0,-1.1){$I$};\end{scope}}
\end{exercise}
\begin{solution}
图\ref{fig:o4mixtureq}两圆态各有$S_1=S_2=0$，$S_3$分别等于其强度的正、负值。直接加斯托克斯量，得$S_0=4I,S_1=S_2=0,S_3=2I$，偏振度$1/2$。
\Diagram{fig:o4mixturea}{线检偏透过恒为$2I$，波片分析后才显出圆偏振不平衡}{\draw[axis](0,0)--(4.6,0)node[right]{检偏角};\draw[axis](0,0)--(0,2.8)node[above]{透射强度};\draw[very thick](0,1.7)--(4,1.7);\node[left]at(0,1.7){$2I$};}
线检偏强度恒为$S_0/2=2I$，如图\ref{fig:o4mixturea}。虽然它的旋转曲线与自然光相同，光束仍有一半的圆偏振成分；加四分之一波片后两线方向透过强度成为$3I$与$I$，即可识别。独立光束按斯托克斯量相加；Jones矢量直接相加只适用于有确定相位关系的场。
\end{solution}

\begin{exercise}\label{ex:o4stack}
自然光从空气入射，穿过一叠相互平行的无损玻璃片后回到空气，每个界面均处于相应方向的布儒斯特角。设共经历$M$个界面（$M$为偶数），忽略反射光重新进入透射主束。若每界面$s$偏振反射率为$R_s$，求直接透射主束光强和线偏振度；另说明若在玻璃内部截取主束，光强应怎样换算。
\Diagram{fig:o4stackq}{玻片堆逐渐削弱$s$分量，保留$p$分量}{\foreach\x in{0,1,2,3}{\draw[curve](\x,-.9)--({\x+.5},.9);}\draw[axis](-1,0)--(4.5,0)node[right]{主束};\node[above]at(1.8,1.3){每面为布儒斯特角};}
\end{exercise}
\begin{solution}
图\ref{fig:o4stackq}的自然光初始$p,s$各占$I_0/2$。每界面的法向能流透射率为$T_p=1$、$T_s=1-R_s$，而沿光束方向定义的强度还须换算截面。若第$j$面前、后的光线与共同法线分别成$\theta_{j-1},\theta_j$，有
\[
 I_{\sigma,j}\cos\theta_j
 =T_{\sigma,j}I_{\sigma,j-1}\cos\theta_{j-1},
 \qquad \sigma=p,s.
\]
逐面相乘，所有中间余弦因子抵消。最终又回到空气，$\theta_M=\theta_0$，所以
\[
 I_p=\frac{I_0}{2},\qquad
 I_s=\frac{I_0}{2}(1-R_s)^M.
\]
\Diagram{fig:o4stacka}{两正交强度份额决定偏振度}{\draw[curve](0,0)rectangle(4,.6);\node at(2,.3){$I_p=I_0/2$};\draw[curve](0,-1)rectangle(1.5,-.4);\node[right]at(1.7,-.7){$I_s=(I_0/2)(1-R_s)^M$};}
总强度$I=I_0[1+(1-R_s)^M]/2$，线偏振度$\mathcal P=[1-(1-R_s)^M]/[1+(1-R_s)^M]$。图\ref{fig:o4stacka}在$0<R_s<1$、$M\to\infty$时趋于完全$p$偏振、主束强度$I_0/2$。

若在玻璃内的第$j$面后截取主束，上述两个强度须共同乘$\cos\theta_0/\cos\theta_j$，并将幂次$M$改为$j$；偏振度中共同因子抵消。这样便保留了法向通量与光束强度的区别。若多次反射旁束也进入探测孔径，须重新求和；有吸收或表面误差时，$p$分量也会损失。
\end{solution}

\begin{exercise}\label{ex:o4partialmeasure}
未知光由自然光和某方向线偏振光组成。旋转检偏器得到最大强度8、最小强度2，单位相同。求入射总强度、自然成分和线偏振成分。若去掉“线偏振”这一前提，还能唯一确定偏振度吗？
\Diagram{fig:o4partialq}{检偏器扫描得到非零的最低强度}{\draw[axis](0,0)--(5.2,0)node[right]{角度};\draw[axis](0,0)--(0,3)node[above]{透射强度};\draw[curve,domain=0:180,samples=80]plot({\x/40},{(2+6*cos(\x)^2)/3});\node[left]at(0,2.667){$8$};\node[left]at(0,.667){$2$};}
\end{exercise}
\begin{solution}
按图\ref{fig:o4partialq}，$I_u/2=2$，所以$I_u=4$；$I_p+I_u/2=8$给$I_p=6$。入射总强度$I_0=10$，偏振度$6/10=.6$。
\Diagram{fig:o4partiala}{线检偏扫描测得$S_1,S_2$，不能单独得到$S_3$}{\draw[curve](0,0)rectangle(2.5,1);\node at(1.25,.5){线检偏器};\draw[->](2.6,.5)--(3.5,.5);\node[right]at(3.5,.5){$S_0,S_1,S_2$};\node[below]at(2,-.6){圆偏振信息需波片};}
若完全偏振分量允许为椭圆或圆偏振，图\ref{fig:o4partiala}说明圆偏振信息$S_3$未测得。单次扫描只能给线性调制度.6，即完整偏振度的下限；它可以介于.6与1之间。附加适当四分之一波片测圆偏振分析强度才能完成判定。
\end{solution}

\begin{exercise}\label{ex:o4plateorder}
线偏振沿$x$，先过快轴45度的四分之一波片，再过快轴0度的半波片。求输出Jones矢量，并比较交换两片顺序的结果。
\Diagram{fig:o4orderq}{两块转角不同的波片，矩阵按传播次序从右向左相乘}{\draw[axis](-1,0)--(5.5,0);\draw[curve](0,-1)--(0,1);\draw[curve](2.8,-1)--(2.8,1);\node[above]at(0,1){$Q(45^\circ)$};\node[above]at(2.8,1){$H(0)$};\node[left]at(-1,.5){$x$线偏振};}
\end{exercise}
\begin{solution}
由旋转矩阵，$Q(45^\circ)=\tfrac12\bigl(\begin{smallmatrix}1+i&1-i\\1-i&1+i\end{smallmatrix}\bigr)$，$H(0)=\operatorname{diag}(1,-1)$。图\ref{fig:o4orderq}的输出为$\tfrac12(1+i,-1+i)^{\mathsf T}$，删去共同相位可写$(1,i)^{\mathsf T}/\sqrt2$。
\Diagram{fig:o4ordera}{交换顺序使最终圆偏振的旋向相反}{\begin{scope}[xshift=-2cm]\draw[curve](0,0)circle(.9);\draw[->](.9,0)arc(0:70:.9);\node[below]at(0,-1.2){先$Q$后$H$};\end{scope}\begin{scope}[xshift=2cm]\draw[curve](0,0)circle(.9);\draw[->](.9,0)arc(0:-70:.9);\node[below]at(0,-1.2){先$H$后$Q$};\end{scope}}
交换后$H(0)$先作用于$(1,0)$不改偏振，输出为$\tfrac12(1+i,1-i)^{\mathsf T}$，删共同相位为$(1,-i)^{\mathsf T}/\sqrt2$。图\ref{fig:o4ordera}的两箭头相反；不对易来自两波片的本征轴不同，而不是共同相位差。
\end{solution}

\begin{exercise}\label{ex:o4normalenergy}
空气向折射率$n=1.5$玻璃正常入射。求振幅反射、透射系数以及能量反射、透射率。再解释为何透射场振幅小于入射场并不表示透过能量只有振幅比的平方。
\Diagram{fig:o4normalq}{正常入射界面，采用共同切向电场基底}{\draw[curve](2,-1.4)--(2,1.4);\draw[->](0,.6)--(1.9,.6);\draw[->](1.9,0)--(0,0);\draw[->](2.1,.6)--(4.5,.6);\node[above]at(.8,1){$n_1=1$};\node[above]at(3.3,1){$n_2=1.5$};}
\end{exercise}
\begin{solution}
图\ref{fig:o4normalq}统一切向场有$r=(1-1.5)/(1+1.5)=-.2$、$t=2/(1+1.5)=.8$。所以$R=.04$，$T=(n_2/n_1)t^2=1.5\times.64=.96$。
\Diagram{fig:o4normala}{能量份额为4\%反射、96\%透射}{\draw[curve](0,0)rectangle(5,1);\draw[curve](.2,0)--(.2,1);\node[above]at(.1,1){4\%};\node at(2.6,.5){96\%};\node[below]at(2.5,0){$R+T=1$};}
图\ref{fig:o4normala}满足能量守恒。平面波光强为$n\varepsilon_0cE_0^2/2$，介质2具有不同波导纳，因此$t^2=.64$只是电场振幅平方比。若使用局部$p$基，振幅$r_p=+.2$，能量反射率仍相同，必须同时注明基向量改变。
\end{solution}

\begin{exercise}\label{ex:o4brewsternumerical}
自然光从空气射向$n=1.5$玻璃，入射角为布儒斯特角。求该角、折射角、反射光强相对入射光强的比例及反射偏振方向。
\Diagram{fig:o4brewsterq}{布儒斯特条件使反射与透射光线垂直}{\draw[curve](-2.5,0)--(2.5,0);\draw[guide](0,-1.7)--(0,2);\draw[->](-2,-1.333)--(0,0);\draw[->](0,0)--(2,-1.333);\draw[->](0,0)--(1.2,1.8);\node[right]at(2,-1.333){反射};\node[right]at(1.2,1.8){透射};}
\end{exercise}
\begin{solution}
图\ref{fig:o4brewsterq}满足$\tan i_B=1.5$，所以$i_B\simeq56.31^\circ$、$r_B\simeq33.69^\circ$。$p$反射为零，$s$振幅系数代入$\cos i_B=1/\sqrt{1+n^2}$、$\cos r_B=n/\sqrt{1+n^2}$得$r_s=(1-n^2)/(1+n^2)=-5/13$。
\Diagram{fig:o4brewstera}{自然光的两个偏振份额中仅$s$有反射}{\draw[curve](0,0)rectangle(2.8,.7);\node at(1.4,.35){$I_0/2$的$s$};\draw[->](2.9,.35)--(4,.35);\node[right]at(4,.35){乘$25/169$};\node[below]at(2,-.7){$p$反射为0};}
自然光的$s$份额只有$I_0/2$，故总反射率为$(1/2)(25/169)=25/338\simeq.07396$，如图\ref{fig:o4brewstera}。反射光完全为$s$偏振，电场垂直入射面；不能把$R_s$直接当自然光的总反射率。
\end{solution}

\begin{exercise}\label{ex:o4retardthickness}
某波片在目标波长$\lambda_0$处双折射差为$\Delta n$。求零级四分之一波片厚度及厚度误差$\Delta h$引起的相位误差。若忽略材料色散，偏离波长时会怎样？
\Diagram{fig:o4thicknessq}{波片中快慢分量走过相同厚度却积累不同相位}{\draw[curve](1,-1)rectangle(3,1);\draw[->](0,.4)--(4.5,.4)node[right]{快};\draw[->](0,-.4)--(4.5,-.4)node[right]{慢};\draw[<->](1,-1.4)--(3,-1.4)node[midway,below]{$h$};}
\end{exercise}
\begin{solution}
图\ref{fig:o4thicknessq}中令$\Delta n=n_s-n_f>0$，慢轴相对快轴的相位延迟为$\delta=2\pi\Delta n h/\lambda_0$。零级四分之一波片取最小正延迟$\pi/2$，故$h_0=\lambda_0/(4\Delta n)$。固定波长、材料时厚度误差给
\[
 \Delta\delta=\frac{2\pi\Delta n}{\lambda_0}\Delta h
 =\frac\pi2\frac{\Delta h}{h_0}.
\]
若快慢轴命名相反，延迟符号随之改变，最小厚度应使用$|\Delta n|$。
\Diagram{fig:o4thicknessa}{延迟偏离90度后，45度线偏振不再变成严格圆偏振}{\draw[axis](-1.5,0)--(1.7,0);\draw[axis](0,-1.5)--(0,1.7);\draw[guide](0,0)circle(1);\draw[curve,rotate=45](0,0)ellipse(1.2 and .75);\node[right]at(1.8,.7){圆变椭圆};}
图\ref{fig:o4thicknessa}中，45度线偏振入射后的Jones矢量正比$(1,\ee^{i\delta})$。若$\delta=\pi/2+\epsilon$，其归一化$S_3/S_0=\sin\delta\simeq1-\epsilon^2/2$，$S_2/S_0=\cos\delta\simeq-\epsilon$，不再严格为圆偏振。

固定厚度且忽略色散时，$\delta(\lambda)=(\pi/2)\lambda_0/\lambda$，从$1/(\lambda_0+\Delta\lambda)$一阶展开得$\Delta\delta\simeq-(\pi/2)\Delta\lambda/\lambda_0$。保留材料色散则应微分整个比值：
\begin{align*}
 \frac{\dd\delta}{\dd\lambda}
 &=2\pi h\left[\frac1\lambda\frac{\dd\Delta n}{\dd\lambda}-\frac{\Delta n}{\lambda^2}\right]\\
 &=\delta\left[\frac1{\Delta n}\frac{\dd\Delta n}{\dd\lambda}-\frac1\lambda\right].
\end{align*}
多级波片的设计延迟为$\pi/2+2m\pi$，其相对厚度误差或相对波长偏移对应更大的绝对相位误差。容差由完整延迟的导数确定；先模去$2\pi$会丢失多级波片的灵敏度。
\end{solution}

